\chapter{双臂与多臂协同操作}
\label{ch:mr-multiarm}

% ════════════════════════════════════════════════════════════════
%  全吸收章：重渲染 Tutorial 笔记
%   - 04_移动机器人规控/50_多机器人协作/80_双臂多臂协同操作.md
%     （视角重置 / 对称表述 / 相对雅可比 / ker(J_rel)-ker(G) 对偶 / 内力控制 /
%      主从-协调架构 / 对称-非对称 / N 臂 QP+ADMM 分配 / 学习范式简述）
%  记号归一：R∈SO(3)、T∈SE(3)（preamble 宏 \SO \SE \so \se）；
%            twist 顺序 [v;ω]（线在前、角在后，源文约定，保留并就地声明）；
%            机器人/操作领域自然记号保留。
%  跨章去重：双臂工程实现（Pinocchio C++/奇异/ROS2）属机械臂方向，本章不重复，
%            只在该处 \pz 指向；共识/ADMM/grasp matrix 的最深推导留待对应章。
%  源由 AI 撰写，作者归属/年份/卷期/数值/前沿编号存疑处就地 \pz/\rebuilt，并入 claims 台账。
%  教学层遵循 v5.0 理论教学规范（动机→反面→历史→理论→陷阱→练习 + 误解汇总/小结/速查/故障排查）。
% ════════════════════════════════════════════════════════════════

\rebuilt{本章重渲染自 Tutorial 笔记《双臂与多臂协同操作》，其源稿由 AI 撰写，论文归属、年份、卷期与前沿工作的编号可能有误。本章忠实重述、不擅自“修正”，逐处 \texttt{\textbackslash cite} 标源；存疑断言已登记 claims 台账，待联网核对后二次校验。本章站在\emph{多机协作}立场重审“双臂操作”，把它当作 $N=2$ 的多机协同操作；与“把双臂当一台 14 自由度机器”的单系统工程视角（机械臂方向）互补而非重复。}

\begin{practice}[前置自测]
开始之前，先试答下列问题。若已能流畅作答，可只速读本章表格与速查卡；若卡住（答不出 $\ge 2$ 题），第 1、3 题回协同搬运章（grasp matrix），第 2 题回共识章（\cref{ch:mr-consensus}），第 4、5 题回机械臂方向的协调运动学/力控基础，再来读本章。本章会大量复用这些前置，但不重新教它们。
\begin{enumerate}
  \item 写出 grasp matrix $\mathbf{G}$ 的定义：$N$ 个接触点的接触力 $\mathbf{f}_{\text{c}}\in\mathbb{R}^{3N}$ 如何映射到负载 wrench $\mathbf{F}_{\text{load}}\in\mathbb{R}^6$？零空间 $\ker(\mathbf{G})$ 的维度是多少、物理上代表什么？
  \item 什么是\textbf{最大值共识}？它与平均共识的收敛目标有何区别？为什么“在通信图上传播一个全局极值”所需轮数正比于图的\textbf{直径}？
  \item 内力/运动力分解 $\mathbf{f}_{\text{c}}=\mathbf{G}^{+}\mathbf{F}_{\text{load}}+(\mathbf{I}-\mathbf{G}^{+}\mathbf{G})\mathbf{f}_{\text{int}}$ 中，为什么 $(\mathbf{I}-\mathbf{G}^{+}\mathbf{G})\mathbf{f}_{\text{int}}$ 这一项“不产生负载运动”？把它代回 $\mathbf{F}_{\text{load}}=\mathbf{G}\mathbf{f}_{\text{c}}$ 验证。
  \item 双 7 自由度臂共持一个刚性物体，约束流形维度是多少？相对雅可比 $\mathbf{J}_{\text{rel}}$ 的零空间 $\ker(\mathbf{J}_{\text{rel}})$ 在物理上代表什么关节运动？
  \item 阻抗控制与导纳控制的二分法是什么？两臂之间夹着物体时，物体扮演“环境”角色——这时应该用阻抗还是导纳？为什么？
\end{enumerate}
\noindent\emph{答案线索}：1$\to$\cref{sec:ma-perspective,sec:ma-duality}（grasp matrix 与 $\ker(\mathbf{G})$ 维度 $3N-6$）；2$\to$\cref{sec:ma-arch}（最大值共识与信息传播）；3$\to$\cref{sec:ma-duality}（伪逆性质 $\mathbf{G}\mathbf{G}^{+}\mathbf{G}=\mathbf{G}$）；4$\to$\cref{sec:ma-coopkin}（$14-6=8$ 维约束流形）；5$\to$\cref{sec:ma-internalforce}（应用阻抗）。
\end{practice}

\section*{本章目标}
学完本章，你应当能够：
\begin{enumerate}
  \item \textbf{用多机协作的语言重述双臂操作}：把“两臂共持物体”形式化为“$N=2$ 个智能体通过共享刚性负载耦合”的特例，说明 grasp matrix $\mathbf{G}$ 与相对雅可比 $\mathbf{J}_{\text{rel}}$ 是同一闭链约束的两种投影——前者在\emph{力空间}、后者在\emph{速度空间}，并经运动学-静力学对偶相连；
  \item \textbf{推导相对雅可比与绝对/相对分解}：从 Uchiyama--Dauchez 对称表述出发，导出绝对变量（物体整体运动）与相对变量（两臂相对运动）的对偶分解，给出 Jamisola--Roberts 紧凑相对雅可比的来历，并理解非对称任务下的扩展相对雅可比；
  \item \textbf{统一闭链约束的两个面孔}：用虚功原理证明 $\ker(\mathbf{J}_{\text{rel}})$（速度零空间）与 $\ker(\mathbf{G})$（内力空间）互为对偶，说明“刚体把持”为什么同时是 $6$ 维速度约束与 $6$ 维 wrench 约束；
  \item \textbf{设计内力控制律并推广到 $N$ 臂}：推导双臂内力 $\mathbf{f}_{\text{int}}=\tfrac12(\mathbf{h}_R-\mathbf{h}_L)$ 的物理含义，给出 Caccavale 内/外阻抗双层结构，并把它推广为 $N$ 臂内力调节，理解为什么内力不能开环施加、必须闭环维护；
  \item \textbf{在“主从 vs 协调”的架构辩论中做出选型}：用集中/分布/去中心化三分法分析双臂/多臂控制，说明主从何时够用、何时因“从臂只跟随不协商”而积累内力，对称协调如何用绝对/相对解耦消除此问题，leader--follower 一致性如何把主从“软化”成可扩展的分布式方案；
  \item \textbf{解决多臂（$N\ge 3$）负载分配问题}：把“$N$ 只手怎么分担负载力”形式化为 grasp 伪逆零空间参数化 $+$ QP 优化，并理解分布式版本（ADMM/共识）与协同搬运、分布式优化的接合点；
  \item \textbf{厘清本章与单系统双臂内容、与协同搬运/多足/人形章节的边界}，理解“双臂 $\to$ 多臂 $\to$ 多移动机械臂 $\to$ 多四足搬运”是同一套闭链数学在不同硬件上的实例化。
\end{enumerate}

\subsection*{本章知识导航}
\textbf{一句话定位}：本章把“协同”从“一群机器人各自运动、互不接触”推进到\emph{它们通过一个共享的刚性负载在物理上锁死在一起}——这是多机耦合最强的一种形式，强到可以写出精确的代数约束（闭链）。双臂是这种耦合的“最小完整实例”：$N=2$ 已经包含内力、冗余、协调、主从全部核心矛盾，却小到能在一页纸上把矩阵写全。

\begin{center}
\small
\begin{tikzpicture}[
  node distance=5mm and 7mm, >=Stealth,
  blk/.style={draw, rounded corners, align=center, font=\small, inner sep=3pt, fill=main!6, draw=main!60!black},
  grp/.style={draw, rounded corners, align=center, font=\small, inner sep=3pt, fill=second!6, draw=second!60!black},
  fin/.style={draw, rounded corners, align=center, font=\small, inner sep=3pt, fill=third!8, draw=third!60!black},
  ln/.style={->, thick, gray!60!black}]
\node[blk] (persp) {视角重置\\ 双臂 $=N{=}2$ 多机协同};
\node[grp, below=8mm of persp] (kin) {协同运动学\\ 对称表述 $\to J_{\text{abs}},J_{\text{rel}}$};
\node[grp, right=of kin] (dual) {约束两面孔\\ $\ker(J_{\text{rel}})\leftrightarrow\ker(G)$};
\node[grp, below=8mm of kin] (force) {内力控制\\ Caccavale 内/外阻抗};
\node[fin, right=of force] (arch) {架构辩论\\ 主从 / 一致性 / 协调};
\node[fin, below=8mm of force] (asym) {对称/非对称\\ $\alpha$ 统一};
\node[fin, right=of asym] (alloc) {$N$ 臂分配\\ QP $+$ ADMM};
\draw[ln] (persp)--(kin); \draw[ln] (kin)--(dual);
\draw[ln] (kin)--(force); \draw[ln] (dual.south)--(arch.north);
\draw[ln] (force)--(arch); \draw[ln] (force)--(asym);
\draw[ln] (asym)--(alloc); \draw[ln] (arch.south)--(alloc.north);
\end{tikzpicture}
\end{center}

\noindent\textbf{推荐路径}：\cref{sec:ma-perspective} 是视角地基（必读，它决定你怎么“看”后面所有内容）；\cref{sec:ma-coopkin} 是数学引擎（必读，相对雅可比是全章核心对象）；\cref{sec:ma-duality} 把力视角与速度视角缝起来（理论方向必读，工程方向可速读结论）；\cref{sec:ma-internalforce} 内力控制是“会用”的核心；\cref{sec:ma-arch} 主从 vs 协调是体系结构思维的训练（与图论/共识章呼应）；\cref{sec:ma-symmetry} 对称/非对称在做真实任务（拧瓶盖、装配）时必读；\cref{sec:ma-allocation} 是推广到 $N\ge3$ 的收口；\cref{sec:ma-boundary} 厘清边界与学习范式。\cref{sec:ma-coopkin,sec:ma-internalforce,sec:ma-arch,sec:ma-allocation} 是四个硬核小节，值得花最多时间。

\subsection*{前置知识桥接}
本章紧接协同搬运（grasp matrix）与机械臂方向的协调运动学，这里把最关键的前置点重新激活——你不必翻回去就能跟上；重叠处一律交叉引用指过去、不重复：
\begin{itemize}
  \item \textbf{grasp matrix 与内力/运动力分解}：协同搬运章用 $\mathbf{G}\in\mathbb{R}^{6\times 3N}$ 把 $N$ 个接触点的力映射到负载 wrench $\mathbf{F}_{\text{load}}=\mathbf{G}\mathbf{f}_{\text{c}}$，并把接触力分解为运动力 $\mathbf{G}^{+}\mathbf{F}_{\text{load}}$ 与内力 $(\mathbf{I}-\mathbf{G}^{+}\mathbf{G})\mathbf{f}_{\text{int}}\in\ker(\mathbf{G})$。本章把双臂当成 $N=2$ 特例：两只手是两个接触点，$\mathbf{G}\in\mathbb{R}^{6\times 12}$（每接触是 $6$ 维 wrench，因夹爪能传力矩），内力空间维度 $12-6=6$，恰对应两臂间 $6$ 维相对位姿可被“较劲”。本章把它和\emph{速度空间}（相对雅可比）对偶起来。\pz{协同搬运章（grasp matrix、内/运动力分解、自适应分配）的最深推导待该章写就后改交叉引用指过去。}
  \item \textbf{共识作为“分布式裁决”的原子操作}（\cref{ch:mr-consensus}）：共识让各 agent 只靠邻居通信收敛到一致（平均共识 $\dot{\mathbf{x}}=-\mathbf{L}\mathbf{x}$、最大值共识）。本章 \cref{sec:ma-arch} 的 leader--follower 内力调节本质是\emph{带物理耦合的一致性}：每个 agent 维护“自己认为物体该怎么动”的估计，通过共享负载的接触力（隐式通信）或显式邻居通信不断刷新，最终对“物体运动 $+$ 内力分配”达成一致——“边”既可以是通信链路，也可以是穿过刚体的力。
  \item \textbf{分布式优化（ADMM）}（\cref{ch:mr-consensus}）：ADMM 把一个大的联合优化拆成各 agent 的本地子问题 $+$ 一致性约束（\cref{sec:cs-admm}）。本章 \cref{sec:ma-allocation} 的多臂负载分配，集中式版本是一个 QP（在 grasp 零空间里优化内力满足摩擦锥），分布式版本正是用 ADMM/对偶分解把该 QP 拆给各 agent。
  \item \textbf{相对雅可比与零空间投影}（机械臂方向）：单系统视角已给出 $\mathbf{J}_{\text{rel}}=\bar{\mathbf{J}}_R-\bar{\mathbf{J}}_L$（先用平移算子把两臂 twist 统一到公共参考点，再做差）与零空间投影 $\mathbf{P}=\mathbf{I}-\mathbf{J}_{\text{rel}}^{+}\mathbf{J}_{\text{rel}}$。本章从\emph{对称表述}重新推导它，让“$\mathbf{J}_{\text{rel}}=\bar{\mathbf{J}}_R-\bar{\mathbf{J}}_L$”成为 Uchiyama--Dauchez 对称框架的一个特例，并推广到非对称任务与 $N$ 臂。\pz{机械臂方向的双臂工程实现（$\mathbf{J}_{\text{aug}}$ 的 Pinocchio C++、奇异/阻尼伪逆、约束规划、ROS2 力控）属单系统深度，本章不重复，待该方向章节写就后改交叉引用。}
\end{itemize}

\subsection*{如果跳过本章会怎样}
\begin{itemize}
  \item \textbf{把双臂当两个独立单臂，物体被“撕裂”}：让人形双手端托盘，左右手各跑一个独立笛卡尔阻抗、互不通气。结果两个控制器各自的微小跟踪误差在\emph{相对位姿}上累积，几秒后托盘被一拉一推——内力失控增长，物体被挤碎或滑脱。\cref{sec:ma-coopkin,sec:ma-internalforce} 会告诉你：必须把关节当成整体，用相对雅可比把运动拆成“物体运动（绝对）$+$ 相对运动（内力）”两个解耦通道，对前者做运动控制、对后者做力控——这正是单臂经验\emph{不能}直接照搬到双臂的地方。
  \item \textbf{多臂搬运时力分配拍脑袋，某只手率先饱和}：四台移动机械臂抬大钢板，“四只手平均分担、每只手 $1/4$”。但钢板质心偏置、各臂能力不同，平均分配让最不利的那只手最先力矩饱和，份额瞬间转嫁邻居引发连锁过载。\cref{sec:ma-allocation} 会告诉你：$N\ge3$ 时负载分配有 $3N-6$（含力矩 $6N-6$）维内力自由度可优化，应用 QP 在 grasp 零空间里最小化“最大归一化接触力”，而非均分；要在没有中央节点时也能做，得用 ADMM 把 QP 分布式化。
\end{itemize}

\begin{table}[H]\centering\small
\caption{本章阅读方式建议}
\label{tab:ma-reading}
\begin{tabular}{lll}
\toprule
阅读方式 & 时间 & 适合谁 \\
\midrule
精读（跟着推完对称分解、$\ker(\mathbf{J}_{\text{rel}})/\ker(\mathbf{G})$ 对偶、内力控制律、$N$ 臂 QP/ADMM） & 18--26 小时 & 第一次系统学协同操作的算法工程师 \\
速读（读懂相对雅可比、绝对/相对分解、主从 vs 协调取舍、$N$ 臂 QP 形式） & 5--7 小时 & 有机械臂背景、补“多机视角” \\
速查（相对雅可比 $+$ 内力分解 $+$ 主从/协调对比表 $+$ $N$ 臂 QP 模板 $+$ 关系表 $+$ 故障排查） & 40--60 分钟 & 实现协同力控时回查 \\
\bottomrule
\end{tabular}
\end{table}

\begin{note}
\textbf{本章记号与约定.}\;
旋转矩阵记 $\mathbf{R}\in\SO(3)$、位姿 $\mathbf{T}\in\SE(3)$（preamble 宏 $\SO,\SE,\so,\se$，与全书一致，见 \nameref{ch:notation}）；机器人/操作领域的自然记号保留。
\textbf{twist 顺序}：本章 twist $\boldsymbol{\mathcal{V}}=[\mathbf{v};\boldsymbol{\omega}]$ 取\emph{线速度在前、角速度在后}（沿用源文与协同操作文献惯例；这与全书 $\se(3)$ 坐标 $\boldsymbol{\xi}=[\boldsymbol{\rho};\boldsymbol{\phi}]$ 的“平移在前”一致，但请注意这里是\emph{速度} twist 而非 $\se(3)$ 指数坐标，$\boldsymbol{\rho}\neq\mathbf{t}$ 的区分仍成立）。若改用 $[\boldsymbol{\omega};\mathbf{v}]$ 顺序，下文平移算子与伴随矩阵的分块位置须相应置换。
臂/接触点数记 $N$（双臂 $N=2$），下标 $i$；左右臂记 $L,R$。反对称算子 $(\cdot)^\wedge$（$\mathbf{a}\times\mathbf{b}=\mathbf{a}^\wedge\mathbf{b}$），$[\mathbf{d}]_\times:=\mathbf{d}^\wedge$。
\textbf{字母隔离声明}：本章 $\alpha$ 指绝对变量的加权参数（非 class-$\mathcal{K}$ 系数）、$\boldsymbol{\lambda}$ 在 \cref{sec:ma-allocation} 指内力参数或 ADMM 对偶变量（按上下文，就地申明）、$\mathbf{G}$ 指 grasp matrix（非通信图，通信图记 $\mathcal{G}$）；均仅本章局部生效。
\end{note}

% ════════════════════════════════════════════════════════════════
\section{视角重置：把双臂当成 $N=2$ 的多机协同操作}
\label{sec:ma-perspective}

\paragraph{动机：同一系统的两种看法导出不同工程。}
同一个“两臂共持物体”的系统，可以有两种看法：(A) 它是\emph{一台 $14$ 自由度的机器}，用一套大雅可比、大动力学描述（单系统工程视角）；(B) 它是\emph{两个智能体通过共享负载耦合的多机系统}（本章立场）。设想两台 $7$ 自由度臂端着一块玻璃，问三个问题：玻璃往哪走、走多快？两只手“夹”玻璃的力多大？若其中一只手电机临时过热降功率，系统怎么办？视角 (A) 会把“玻璃位姿”和“夹持力”塞进同一个任务向量求解——前两问没问题，但第三问别扭：一台机器的“一个电机降功率”只是某关节力矩约束变了，得重组整个 QP，它没有“某 agent 能力下降、把份额转给邻居”这种\emph{天然的模块边界}。

\paragraph{多机视角下三个问题对应三个熟悉框架。}
持视角 (B)，三问立刻对应到已学的多机工具（\cref{tab:ma-translate-q}）。这不是文字游戏：视角决定模块边界，模块边界决定代码结构、容错策略与可扩展性。

\begin{table}[H]\centering\small
\caption{双臂三问题的多机协作对应}
\label{tab:ma-translate-q}
\begin{tabular}{lll}
\toprule
问题 & 多机协作视角下的对应 & 出处 \\
\midrule
玻璃往哪走、多快 & 团队对“共同目标（物体轨迹）”达成一致 & 共识 / 分布式 MPC（\cref{ch:mr-consensus}） \\
两手夹持力多大 & 内力分配 $=$ grasp 零空间里的力优化 & grasp matrix $+$ \cref{sec:ma-internalforce,sec:ma-allocation} \\
一只手降功率怎么办 & 某 agent 能力下降 $\to$ 重新分配负载份额 & 自适应分配 $+$ \cref{sec:ma-allocation} QP \\
\bottomrule
\end{tabular}
\end{table}

\begin{insight}[双臂是多机耦合的最小完整实例]
双臂操作之所以是“多机协作”的最小完整实例，是因为它在 $N=2$、自由度最少、还能把所有矩阵写满的前提下，\emph{同时}具备多机系统的全部核心矛盾——共享目标的一致（物体往哪走）、资源的分配（力怎么分）、容错与异构（一只手坏了 / 能力不同）。它比“两台机器人在走廊里避让”耦合强得多（那里只有碰撞这种不等式弱耦合，这里是穿过刚体的等式强耦合），又比“十台机器人搬一栋楼”简单得多（$N=2$ 没有可扩展性爆炸）。研究双臂，是在最干净的实验台上研究多机耦合本身。
\end{insight}

\paragraph{反面：只用视角 (A) 会看不见三类结构。}
视角 (A) 不是错的，而是\emph{会让你看不见某些结构}，从而在三类情况下吃亏。其一，\textbf{异构与容错被埋没}：两条臂型号不同（臂展、载荷、控制频率不一）、或系统从“两臂”扩展到“两台移动双臂机器人各出一只手”时，视角 (A) 的“一台 $14$ 自由度机器”假设崩塌——它默认所有自由度共享同一个时钟、求解器、全局状态；视角 (B) 从一开始就假设“每个 agent 有自己的控制器、只能观测局部 $+$ 邻居”。其二，\textbf{通信约束被忽略}：单台机器关节间走板内总线，延迟可忽略；真实多臂系统可能挂在不同工控机上经以太网/DDS 交换状态，有毫秒级延迟与丢包，视角 (B)（继承图论/共识）把通信图、延迟、拓扑变化当一等公民。其三，\textbf{与课程其余部分割裂}：用视角 (A) 写本章，它会变成机械臂方向的缩写版，和共识、任务分配、分布式 MPC、grasp matrix 没有连线；用视角 (B)，本章成为这些主题在协同操作上的汇流。

\paragraph{一张“翻译表”落实视角 (B)。}
\cref{tab:ma-translate} 把双臂操作的每个对象“翻译”成多机协作标准语言，后面每节兑现其中某一行。

\begin{table}[H]\centering\small
\caption{双臂操作 $\to$ 多机协作的翻译表}
\label{tab:ma-translate}
\begin{tabular}{lll}
\toprule
双臂操作里的对象 & 多机协作里的对应概念 & 本章兑现处 \\
\midrule
左臂、右臂 & 两个智能体（agent） & 全章 \\
共持的刚性物体 & 把两 agent 锁在一起的共享约束载体（强耦合边） & \cref{sec:ma-coopkin,sec:ma-duality} \\
物体的目标轨迹 & 团队的共同目标（需达成一致） & \cref{sec:ma-coopkin,sec:ma-arch} \\
两手相对位姿不变 & 等式耦合约束 $\to\ker(\mathbf{J}_{\text{rel}})$ / $\ker(\mathbf{G})$ & \cref{sec:ma-coopkin,sec:ma-duality} \\
夹持内力 & 内力 $=$ 共享约束的“对偶变量”，落在 grasp 零空间 & \cref{sec:ma-duality,sec:ma-internalforce} \\
主臂主导、从臂跟随 & 集中式 / leader--follower 架构 & \cref{sec:ma-arch} \\
两臂地位对等、协商运动 & 分布式 / 去中心化（对称协调） & \cref{sec:ma-arch,sec:ma-symmetry} \\
一只手扶、一只手操作 & 异构角色分工（任务分配的退化形式） & \cref{sec:ma-symmetry} \\
$N$ 只手抬大物体的力分配 & 多机资源分配（QP / ADMM 分布式） & \cref{sec:ma-allocation} \\
一只手能力下降 & 异构 $+$ 容错（份额再分配） & \cref{sec:ma-allocation} \\
\bottomrule
\end{tabular}
\end{table}

\paragraph{在“耦合强度谱”上的精确位置。}
把多机问题按“agent 之间耦合的强度”排成一条谱（\cref{fig:ma-spectrum}），本章位置一目了然。

\begin{figure}[H]
\centering
\begin{tikzpicture}[font=\small,>=Stealth,node distance=2mm]
  \draw[->,thick,gray!60!black] (0,0) -- (13.4,0);
  \node[anchor=south west,font=\footnotesize] at (0,0.05) {弱耦合};
  \node[anchor=south east,font=\footnotesize] at (13.4,0.05) {强耦合};
  \foreach \x/\lab/\sub in {
     1.5/{通信耦合}/{共识/任务分配},
     5.0/{碰撞耦合}/{MAPF 避让},
     8.5/{力耦合(点)}/{协同搬运·点接触},
     11.9/{力+运动耦合}/{\textbf{协同操作·刚体把持}}}{
     \fill (\x,0) circle (1.6pt);
     \node[anchor=south,align=center,font=\footnotesize] at (\x,0.12) {\lab};
     \node[anchor=north,align=center,font=\scriptsize,text=gray!50!black] at (\x,-0.08) {\sub};
  }
  \node[anchor=north,align=center,font=\scriptsize] at (11.9,-0.7) {【本章】};
\end{tikzpicture}
\caption{多机耦合强度谱。本章（协同操作，刚体把持）位于最强耦合端：共享负载 $+$ 完整臂运动学，还要精调内力与在手运动。}
\label{fig:ma-spectrum}
\end{figure}

\noindent 与 MAPF（碰撞耦合）比，那里是\emph{不等式弱耦合}（$d(\mathbf{q}_i,\mathbf{q}_j)\ge d_{\min}$，别撞就行），可解耦近似；本章是\emph{等式强耦合}（刚体把持 $\mathbf{h}(\mathbf{q})=\mathbf{0}$），一只手动了另一只手必须精确跟随到亚毫米否则内力爆炸。与协同搬运比更细：搬运常把每个 agent 对负载的作用简化为“一个接触点施一个力/wrench”，关心力怎么分；本章保留每只手是\emph{有完整运动学的冗余臂}这一事实，于是除“力怎么分”，还要管“两臂运动学怎么协调（相对雅可比）”“在手重定向怎么做（相对运动）”“冗余零空间怎么用”——这些在点接触抽象下看不见。本章的闭链约束、相对雅可比、内力分配也是多足协同 loco-manipulation（“臂”换“腿”）与人形全身协同（浮动基座）的\emph{数学母板}，在那两类问题里原样复用，只是雅可比里多了浮动基座的列、约束里多了接触维持。

\begin{pitfall}[思维陷阱：认为“多机视角只是换说法，对实现无影响”]
新手想法：“不管叫‘$14$ 自由度机器’还是‘两个 agent’，最后解的都是同一个 QP，视角是哲学问题。”实际上：视角决定你给系统画的\emph{模块边界}，而模块边界决定代码结构、容错与可扩展性。视角 (A) 自然导出一个单体大求解器（所有自由度耦合在一个 KKT 系统里），加一条臂要重写矩阵维度；视角 (B) 自然导出“每 agent 一个本地控制器 $+$ 一致性层”，加一条臂只是多一个同构模块。当项目从“双工业臂”演化到“两台移动双臂机器人协同”（$4$ 只手、$2$ 个计算节点）时，视角 (A) 的代码几乎要重写，视角 (B) 增量扩展即可。视角是架构决策的源头，不是修辞。
\end{pitfall}

\begin{pitfall}[概念误区：认为“协同操作 $=$ 协同搬运”]
新手想法：“协同搬运已经讲了多机搬物体，本章是不是重复？”实际上：协同搬运关注力分配，常把 agent 简化为施力点，适合“四足/移动平台抬大件”这种“末端运动学不重要、力分配重要”的场景；协同操作关注的是“有灵巧末端的臂如何在共持的同时做精细协调”——在手重定向、内力精调、相对运动控制，这些需要完整的相对雅可比，在点接触抽象下无法表达。两者是“力空间”与“力$+$运动空间”的递进：搬运是操作的“只关心力”的子集，操作多了一层运动学协调。
\end{pitfall}

\begin{exercise}
\begin{enumerate}
  \item \textbf{[概念]} 把“两台移动双臂机器人合作把沙发抬过门”用 \cref{tab:ma-translate} 逐行翻译：谁是 agent？共享约束载体是什么？共同目标、内力、主从/对称分别对应什么？该系统在耦合强度谱上落在哪、和纯双臂相比多了哪些维度？
  \item \textbf{[对比]} 列出 $3$ 个“用视角 (A) 更方便”与 $3$ 个“用视角 (B) 更方便”的场景，给出判据（提示：agent 数、是否异构、是否分布式计算、是否需要容错）。
  \item \textbf{[跨章综合]} 论证：双臂的“角色分配（哪只手当主、哪只手当辅）”是不是任务分配问题的一个退化特例（$N=2$ 个 agent、$2$ 个角色）？它和一般任务分配的关键区别是什么（提示：两个角色之间有强物理耦合，不能独立执行）？
\end{enumerate}
\end{exercise}

% ════════════════════════════════════════════════════════════════
\section{科研发展脉络}
\label{sec:ma-history}

在钻进方法前先把研究线的来龙去脉理清——双臂/多臂协同操作有近 $40$ 年历史，知道每个方法“从哪来、解决前人什么痛点、又留下什么”比孤立记名字有用。\cref{tab:ma-history} 从多机协作视角梳理这条线。\rebuilt{下表年份、会议、作者归属转述自源稿，部分待联网核对（已就地 \texttt{\textbackslash pz} 标出存疑项）。}

\begin{table}[H]\centering\small
\caption{双臂/多臂协同操作研究脉络（多机协作视角）}
\label{tab:ma-history}
\begin{tabular}{p{0.9cm}p{4.7cm}p{2.1cm}p{4.6cm}}
\toprule
年份 & 工作 & Venue & 核心贡献 \\
\midrule
1988 & Uchiyama--Dauchez，对称式 hybrid position/force 协调\cite{uchiyama1987symmetric}\pz{存疑: 源稿标 1987, 权威记录为 ICRA 1988(Philadelphia, pp.350--356, DOI 10.1109/ROBOT.1988.12073); 论文/作者/内容归属均确, 仅年份与 bibkey 年号(应为 1988)待订} & ICRA & \textbf{对称表述奠基}：绝对/相对变量分解，明确反对“主从”——两臂地位对等、物体运动与内力解耦 \\
1988 & Khatib，augmented object（object-level 操作空间动力学）\cite{khatib1988object}\rebuilt{原误用键 khatib1987operational(=1987 IEEE J-RA 操作空间表述原文, 不含 augmented object)且误标 IJRR; 已改正确键 khatib1988object(ISRR 1988, MIT Press, pp.137--144); 同主题另有 ACC 1988, 期刊级为 IJRR 1995——非 1988} & ISRR（MIT Press）/ACC & \textbf{augmented object}：把负载作为虚拟物体插入 object-level 动力学，统一多接触点的操作空间动力学 \\
1993 & Uchiyama--Dauchez，对称运动学表述与非主从协调控制\cite{uchiyama1993symmetric} & Adv. Robotics & 把对称表述系统化为完整非主从协调框架——\cref{sec:ma-arch,sec:ma-symmetry} 的理论源头 \\
1996 & Chiacchio 等，多臂系统构型的任务空间动力学分析\cite{chiacchio1996task}\rebuilt{断言张冠李戴, 待修: 此键标题对应的真论文是 IJRR 1991 vol.10(4):410--425(单臂冗余/任务空间增广+任务优先, 非双臂绝对/相对); 而"绝对/相对变量(右减左、增广雅可比)"实出自另一篇 Chiacchio--Chiaverini--Siciliano(无 Sciavicco) 1996 ASME JDSMC 118(4):691--697; 当前键的年/卷/页(1996/15/408--422)亦误。本行内容暂不重建, 待补正确键} & IJRR & \textbf{绝对/相对变量现代形式}：右减左约定、增广雅可比，被工程普遍采用 \\
2008 & Caccavale 等，双臂协作机械臂的 $6$ 自由度阻抗控制\cite{caccavale2008sixdof} & T-Mech & \textbf{内/外阻抗双层}：object impedance $+$ internal impedance，\cref{sec:ma-internalforce} 主线 \\
2015 & Jamisola--Roberts，基于各臂雅可比的紧凑相对雅可比\cite{jamisola2015compact} & RAS & \textbf{紧凑相对雅可比}：直接用各臂自身雅可比拼 $\mathbf{J}_{\text{rel}}$，无需重建整链，工程友好 \\
2015 & Erhart--Hirche，协作多机操作的内力分析与负载分配\cite{erhart2017internal}\pz{存疑: 论文真实, 卷期页(T-RO vol.31 no.5 pp.1238--1243)均确, 但年份应为 2015(DOI 10.1109/TRO.2015.2459412)非 2017; bibkey 年号(宜 erhart2015internal)待订} & T-RO & \textbf{多机内力严格分析}：grasp 伪逆的内力依赖末端运动学，给出无内力广义逆参数化 \\
2019 & Almeida--Karayiannidis 等，扩展相对雅可比的非对称双臂任务\cite{almeida2019asymmetric}\pz{存疑: arXiv:1905.01248 与内容归属确; 但第二作者应为 Karayiannidis 非"Lima", 会议应为 ISRR 2019(pp.18--34)非 IROS, 待订} & ISRR/arXiv & \textbf{扩展相对雅可比}：非对称任务下同时控制相对运动 $+$ 一手绝对运动，\cref{sec:ma-symmetry} \\
2019 & Verginis--Dimarogonas，内力调节的刚性理论视角\cite{verginis2019cooperative}\pz{存疑: 核心结论(grasp 矩阵与刚性矩阵的值空间-零空间关系)确, 但元数据多误: 正确 arXiv 为 1911.01297(源填 1905.01498 是无关论文), 应有第三作者 Zelazo, 期刊为 T-CNS(非 T-CST), 期刊版 2023; "分布式/无需负载惯性"系外推待证} & T-CNS/arXiv & \textbf{刚性理论视角}：grasp 矩阵与刚性矩阵的零空间-值空间关系，分布式内力调节，无需负载惯性 \\
2022 & Zhao--Finn 等，低成本硬件细粒度双臂操作（ALOHA）\cite{zhao2023aloha} & RSS & \textbf{学习范式转向}：低成本遥操作 $+$ ACT，双臂技能从“建模控制”转向“模仿学习” \\
2023-25 & Diffusion Policy / Mobile ALOHA / VLA\cite{chi2023diffusion}\rebuilt{待修(勿原位重建): Diffusion Policy 原文(RSS 2023)实验为单臂(UR5/Franka, 4 任务), "双臂/bimanual"任务仅见于其 IJRR 2024 扩展版; 故把"双臂动作生成"归于该 RSS 2023 文有误, 键/会议本身无误} & RSS/CoRL/arXiv & 移动双臂、VLA 直出双臂动作、leader--follower 动作分解——学习与经典控制融合 \\
\bottomrule
\end{tabular}
\end{table}

\noindent\textbf{两条主线交织}。\emph{经典控制线}从“主从”走向“对称协调”再走向“分布式内力调节”，每一步都在去中心化、减少对单一主导者和全局信息的依赖——这与多机“集中式$\to$分布式$\to$去中心化”的大主题完全同构；\emph{学习线}从 $2022$ 年 ALOHA 起异军突起，把“怎么协调”从人工设计的控制律变成从演示数据里学。本章主体讲经典控制线（它是“可推广、可分析、可移植到多机”的部分），学习线在 \cref{sec:ma-boundary} 简述。\pz{实验室脉络（Tohoku Uchiyama $\to$ Stanford Khatib $\to$ Napoli Siciliano/Caccavale $\to$ TUM Hirche/Erhart $\to$ KTH Dimarogonas/Verginis $\to$ Stanford Zhao/Finn）转述自源稿。核查: 6 节点中 5 节点确(Stanford Khatib、Napoli Siciliano、TUM Hirche/Erhart、KTH Dimarogonas/Verginis、Stanford Zhao/Finn 均属实); Uchiyama 仅 Tohoku 属实, 源稿"早稻田(Waseda)"无据(疑为合著者 Waseda 单位误植)已删。}

% ════════════════════════════════════════════════════════════════
\section{协同运动学：相对雅可比与绝对/相对分解}
\label{sec:ma-coopkin}

\paragraph{动机：独立控制两只手为什么不行。}
最朴素的想法：给左手目标位姿 $\mathbf{T}_L^{\text{des}}(t)$、右手 $\mathbf{T}_R^{\text{des}}(t)$，各跑各的逆运动学。只要两目标“事先算好相容”（对应同一物体位姿），理论上物体就被正确搬运。这在\emph{开环、无误差}的理想世界成立，一接触现实就崩：两个独立控制器各有跟踪误差 $\mathbf{e}_L,\mathbf{e}_R$，而真正伤害物体的不是 $\mathbf{e}_L$ 或 $\mathbf{e}_R$ 本身，而是它们的\emph{差}所对应的相对位姿偏差。设左右手实际位姿为 $\mathbf{T}_L^{\text{des}}\exp(\mathbf{e}_L^\wedge)$、$\mathbf{T}_R^{\text{des}}\exp(\mathbf{e}_R^\wedge)$，亚毫米级跟踪误差（$\sim0.5$\,mm）使两手相对位姿误差约 $\mathbf{e}_R-\mathbf{e}_L$ 可达毫米级；刚性物体被“拉伸/挤压” $1$\,mm，内力 $=$ 刚度 $\times$ 形变，钢的等效刚度可达 $10^7$\,N/m 级，于是内力 $\sim 10^4$\,N——物体被压碎或夹爪打滑。

\begin{insight}[独立控制的根本缺陷：绝对与相对耦合在每只手的误差里]
独立控制把“绝对运动”和“相对运动”\emph{耦合在每只手的位姿误差里}——你无法单独惩罚“相对误差”而放过“绝对误差”。而协同操作的物理本质要求\emph{对这两者用完全不同的控制律}：绝对运动是位置任务（物体要精确到达目标），相对运动是力任务（相对位姿允许微小偏移，但偏移对应的内力要被精确调控）。要分开控制，就必须先有一组把它们解耦的坐标——这正是绝对/相对分解的全部动机。
\end{insight}

\subsection{Uchiyama--Dauchez 对称表述：绝对与相对变量}
\label{sec:ma-symformulation}

1988--1993 年\pz{存疑: 早期对称式 hybrid position/force 方案的权威年份为 ICRA 1988(源稿标 1987), 详见 \cref{tab:ma-history}} Uchiyama 和 Dauchez 提出一个影响至今的思想\cite{uchiyama1987symmetric,uchiyama1993symmetric}：\emph{不要指定一个主、一个从}（主从表述人为打破两臂对称性，\cref{sec:ma-arch} 详述其代价），而是用两个\emph{对称的}物理量描述双臂状态——\textbf{绝对位姿} $\mathbf{T}_{\text{abs}}$（两末端的“中点位姿”，代表被把持物体的整体位姿）与\textbf{相对位姿} $\mathbf{T}_{\text{rel}}$（右末端相对左末端的位姿，刚性把持时应当恒定）。这组坐标“对”，因为它和任务的物理结构对齐（\cref{tab:ma-absrel}）。

\begin{table}[H]\centering\small
\caption{绝对/相对坐标与任务结构的对齐}
\label{tab:ma-absrel}
\begin{tabular}{llll}
\toprule
坐标 & 物理意义 & 任务类型 & 期望行为 \\
\midrule
$\mathbf{T}_{\text{abs}}$（绝对） & 物体往哪走 & 运动任务 & 精确跟踪目标轨迹 \\
$\mathbf{T}_{\text{rel}}$（相对） & 两手“较劲”程度 & 力任务（内力） & 保持恒定（刚体）或受控变化（在手操作） \\
\bottomrule
\end{tabular}
\end{table}

下面给出速度层的精确定义（位姿层用 $\SE(3)$ 对数映射，思路相同但记号繁琐，速度层足以支撑控制）。设两末端 twist 为 $\boldsymbol{\mathcal{V}}_L,\boldsymbol{\mathcal{V}}_R\in\mathbb{R}^6$，\emph{都已通过平移算子统一到公共参考点} $\mathbf{p}_c$（通常取物体质心或两末端中点），记统一后为 $\bar{\boldsymbol{\mathcal{V}}}_L,\bar{\boldsymbol{\mathcal{V}}}_R$。

\begin{note}
\textbf{关键前提：统一参考点.}\;
直接拿 \verb|LOCAL_WORLD_ALIGNED| 雅可比做 $\mathbf{J}_R-\mathbf{J}_L$ 是\emph{错的}——该约定只统一了坐标轴方向，没有把线速度参考点移到同一点。刚体绕物体质心转动时，两末端线速度本来就不同，不统一参考点就做差，会把“正常的绕物体转动”误判成“相对运动”，从而误算出根本不存在的内力。统一参考点的 twist 平移算子（$[\mathbf{v};\boldsymbol{\omega}]$ 顺序）为
\[
  \mathbf{X}(\mathbf{d})=\begin{bmatrix}\mathbf{I}_3 & -[\mathbf{d}]_\times\\ \mathbf{0} & \mathbf{I}_3\end{bmatrix},\qquad
  \mathbf{d}=\mathbf{p}_c-\mathbf{p}_i,\qquad \bar{\boldsymbol{\mathcal{V}}}_i=\mathbf{X}(\mathbf{p}_c-\mathbf{p}_i)\,\boldsymbol{\mathcal{V}}_i.
\]
若用 $[\boldsymbol{\omega};\mathbf{v}]$ 顺序，分块位置与符号要相应置换。
\end{note}

\textbf{绝对 twist（物体整体运动）} 取两末端统一 twist 的平均，\textbf{相对 twist（两手相对运动）} 取右减左：
\begin{equation}
  \boldsymbol{\mathcal{V}}_{\text{abs}}=\tfrac12\big(\bar{\boldsymbol{\mathcal{V}}}_L+\bar{\boldsymbol{\mathcal{V}}}_R\big),
  \qquad
  \boldsymbol{\mathcal{V}}_{\text{rel}}=\bar{\boldsymbol{\mathcal{V}}}_R-\bar{\boldsymbol{\mathcal{V}}}_L.
  \label{eq:ma-absrel-twist}
\end{equation}
这两个定义看似随意，其实有深刻的对称性：它们把 $(\bar{\boldsymbol{\mathcal{V}}}_L,\bar{\boldsymbol{\mathcal{V}}}_R)$ 做了一个\emph{正交变换}（求和取半 $+$ 求差），把“两个末端速度”变成“质心速度 $+$ 相对速度”——和力学里把两质点运动分解为“质心运动 $+$ 相对运动”是同一个数学操作（雅可比坐标变换）。

\begin{insight}[绝对/相对分解 vs 二体问题的质心-相对坐标（带边界）]
在天体力学里，两个相互引力作用的天体不用 $(\mathbf{r}_1,\mathbf{r}_2)$ 而用“质心位置 $\mathbf{R}=\tfrac{m_1\mathbf{r}_1+m_2\mathbf{r}_2}{m_1+m_2}$ $+$ 相对位置 $\mathbf{r}=\mathbf{r}_2-\mathbf{r}_1$”——因为质心做匀速直线运动、相对坐标满足约化质量的单体方程，\emph{两者解耦}。\textbf{像的部分}：都把“两个个体的坐标”线性变换为“集体坐标 $+$ 相对坐标”以解耦。\textbf{不像的部分}：二体问题里两坐标的动力学\emph{严格}解耦（质心绝不受内力影响）；双臂里这种解耦是\emph{控制器刻意制造的}（把关节速度投影到不同子空间），物理上若控制不当，相对运动产生的内力反作用仍会通过物体动力学耦合回绝对运动。\textbf{不要延伸到}：“既然质心运动不受内力影响，物体整体运动也不受内力影响”——那是天体的自然律，双臂里要靠控制器主动维持这个解耦。
\end{insight}

\subsection{从对称变量到雅可比}
\label{sec:ma-jac}

把 \cref{eq:ma-absrel-twist} 写成关于关节速度 $\dot{\mathbf{q}}$ 的线性映射，就得到绝对/相对雅可比。由 $\bar{\boldsymbol{\mathcal{V}}}_i=\bar{\mathbf{J}}_i\dot{\mathbf{q}}$（其中 $\bar{\mathbf{J}}_i=\mathbf{X}(\mathbf{p}_c-\mathbf{p}_i)\mathbf{J}_i$ 是统一参考点后的雅可比，$\mathbf{J}_i\in\mathbb{R}^{6\times n}$ 是第 $i$ 臂在联合关节空间的几何雅可比）：
\begin{align}
  \boldsymbol{\mathcal{V}}_{\text{abs}}&=\tfrac12(\bar{\mathbf{J}}_L+\bar{\mathbf{J}}_R)\dot{\mathbf{q}}=\mathbf{J}_{\text{abs}}\dot{\mathbf{q}},
  & \mathbf{J}_{\text{abs}}&=\tfrac12(\bar{\mathbf{J}}_L+\bar{\mathbf{J}}_R)\in\mathbb{R}^{6\times n},
  \label{eq:ma-Jabs}\\
  \boldsymbol{\mathcal{V}}_{\text{rel}}&=(\bar{\mathbf{J}}_R-\bar{\mathbf{J}}_L)\dot{\mathbf{q}}=\mathbf{J}_{\text{rel}}\dot{\mathbf{q}},
  & \mathbf{J}_{\text{rel}}&=\bar{\mathbf{J}}_R-\bar{\mathbf{J}}_L\in\mathbb{R}^{6\times n}.
  \label{eq:ma-Jrel}
\end{align}
\cref{eq:ma-Jrel} \emph{重新推出了}单系统视角给出的 $\mathbf{J}_{\text{rel}}=\bar{\mathbf{J}}_R-\bar{\mathbf{J}}_L$——但现在你知道它的来历：它是对称表述里“相对变量”的雅可比，“右减左”的符号是相对位姿定义的直接后果，而非随意约定；同时附带得到它的“对偶搭档”$\mathbf{J}_{\text{abs}}$（\cref{eq:ma-Jabs}），这是内力控制必需而单系统推导未强调的。把两者堆叠，就是把 $(\boldsymbol{\mathcal{V}}_{\text{abs}},\boldsymbol{\mathcal{V}}_{\text{rel}})$ 当作新任务坐标：
\begin{equation}
  \begin{bmatrix}\boldsymbol{\mathcal{V}}_{\text{abs}}\\ \boldsymbol{\mathcal{V}}_{\text{rel}}\end{bmatrix}
  =\underbrace{\begin{bmatrix}\tfrac12(\bar{\mathbf{J}}_L+\bar{\mathbf{J}}_R)\\ \bar{\mathbf{J}}_R-\bar{\mathbf{J}}_L\end{bmatrix}}_{\mathbf{J}_{\text{coop}}\in\mathbb{R}^{12\times n}}\dot{\mathbf{q}}.
  \label{eq:ma-Jcoop}
\end{equation}
对双 $7$ 自由度（$n=14$），$\mathbf{J}_{\text{coop}}\in\mathbb{R}^{12\times 14}$，秩 $12$（远离奇异），零空间维度 $14-12=2$——这 $2$ 维是双臂的\emph{冗余}，可用于自碰撞回避、操作度优化、关节限位回避（零空间三大用途的完整推导属机械臂方向，本章不重复）。这个分解把一个 $12$ 维任务拆成两个语义清晰、可分别赋予不同控制律的 $6$ 维通道：绝对通道 $\boldsymbol{\mathcal{V}}_{\text{abs}}$ 做运动控制（跟踪物体目标轨迹），相对通道 $\boldsymbol{\mathcal{V}}_{\text{rel}}$ 做力控（刚体把持时令 $\boldsymbol{\mathcal{V}}_{\text{rel}}=\mathbf{0}$ 保持相对位姿，在手操作时让它跟踪期望相对运动，内力通过相对位姿偏差 $\times$ 刚度产生）。

\subsection{Jamisola--Roberts 紧凑相对雅可比}
\label{sec:ma-compact}

\cref{sec:ma-jac} 用了“先统一参考点 $\bar{\mathbf{J}}_i=\mathbf{X}(\mathbf{p}_c-\mathbf{p}_i)\mathbf{J}_i$ 再做差”。Jamisola 与 Roberts（2015）给出一个\emph{工程上更顺手}的等价表达\cite{jamisola2015compact}：直接用各臂在\emph{各自基座}计算的标准雅可比，配上旋转/wrench 变换矩阵拼出 $\mathbf{J}_{\text{rel}}$，无需把双臂重建成一条长运动链。核心思想：要表达“右末端相对左末端的运动”，需要 (1) 把右臂雅可比变换到左末端参考系；(2) 把左臂雅可比（描述左末端自身运动，而左末端是“相对运动”的参考基座）取负并做相应变换。结果形如
\begin{equation}
  \mathbf{J}_{\text{rel}}=\begin{bmatrix}
  -\,{}^{R_E}\mathbf{T}_{L_E}\,\boldsymbol{\Psi}_L\,{}^{L}\mathbf{J}_L
  & \quad {}^{R_E}\mathbf{T}_{L_E}\,{}^{R}\mathbf{J}_R\end{bmatrix},
  \label{eq:ma-jamisola}
\end{equation}
其中 ${}^{L}\mathbf{J}_L,{}^{R}\mathbf{J}_R$ 是两臂在各自坐标系下的标准几何雅可比，${}^{R_E}\mathbf{T}_{L_E}$ 是把 twist 从左末端系变换到右末端系的 $6\times6$ 伴随（adjoint）矩阵，$\boldsymbol{\Psi}_L$ 是处理“左末端作为相对运动参考基座”时的附加变换。\pz{各变换矩阵的精确形式依赖坐标系约定，细节见 Jamisola--Roberts 2015 原文；本章只取其工程结论。}要记住的不是矩阵的每一项，而是三个工程结论：(1) \textbf{可组合性}——$\mathbf{J}_{\text{rel}}$ 可由两臂“现成的”标准雅可比 $+$ 几个 $\SE(3)$ 变换拼出，不必为双臂单独写一套运动学，代码复用友好；(2) \textbf{左臂列带负号}——左末端是相对运动的“参考者”，其运动以相反方式贡献到相对 twist，这是 $-\bar{\mathbf{J}}_L$ 在紧凑形式里的体现，与第一性原理推导一致；(3) \textbf{任意一臂可作参考}——选左末端、右末端还是物体质心作参考点，只改变变换矩阵，不改变“相对雅可比 $=$ 一臂雅可比变换后减另一臂雅可比变换”的结构。

\begin{note}
\textbf{工程建议（数值真值交叉验证）.}\;
开发期用对称表述 $+$ 显式平移算子（\cref{eq:ma-Jabs,eq:ma-Jrel}）做数值真值，确认 $\mathbf{J}_{\text{rel}},\mathbf{J}_{\text{abs}}$ 正确后，实时控制再切到 \cref{eq:ma-jamisola} 紧凑形式或 C++ 复用单臂雅可比——后者省掉显式平移，但容易在坐标系约定上犯错，必须用前者交叉验证。验证脚手架：构造投影到 $\ker(\mathbf{J}_{\text{rel}})$ 的关节速度（$\mathbf{P}_{\text{rel}}=\mathbf{I}-\mathbf{J}_{\text{rel}}^{+}\mathbf{J}_{\text{rel}}$，取 $\dot{\mathbf{q}}=\mathbf{P}_{\text{rel}}\mathbf{r}$，$\mathbf{r}$ 随机），应有 $\|\mathbf{J}_{\text{rel}}\dot{\mathbf{q}}\|\approx10^{-12}$、$\|\mathbf{J}_{\text{abs}}\dot{\mathbf{q}}\|$ 一般非零。
\end{note}

\subsection{协调逆运动学：两个解耦通道}
\label{sec:ma-coopik}

有了 $\mathbf{J}_{\text{abs}}$ 和 $\mathbf{J}_{\text{rel}}$，刚体把持任务的协调逆运动学就是“绝对通道跟踪物体、相对通道锁零”：
\begin{equation}
  \dot{\mathbf{q}}=\mathbf{J}_{\text{coop}}^{+}\begin{bmatrix}\boldsymbol{\mathcal{V}}_{\text{abs}}^{\text{des}}\\ \boldsymbol{\mathcal{V}}_{\text{rel}}^{\text{des}}\end{bmatrix}
  +\big(\mathbf{I}-\mathbf{J}_{\text{coop}}^{+}\mathbf{J}_{\text{coop}}\big)\dot{\mathbf{q}}_0,
  \label{eq:ma-coopik}
\end{equation}
刚体把持时 $\boldsymbol{\mathcal{V}}_{\text{rel}}^{\text{des}}=\mathbf{0}$，$\dot{\mathbf{q}}_0$ 是冗余自任务（自碰撞回避等）。等价地，可用分层（任务优先级）写法：先在 $\ker(\mathbf{J}_{\text{rel}})$ 里满足相对约束，再在其中实现绝对运动，
\begin{equation}
  \dot{\mathbf{q}}=\underbrace{\mathbf{P}_{\text{rel}}}_{\mathbf{I}-\mathbf{J}_{\text{rel}}^{+}\mathbf{J}_{\text{rel}}}\,\mathbf{J}_{\text{abs}}^{+}\,\boldsymbol{\mathcal{V}}_{\text{abs}}^{\text{des}}+\cdots.
  \label{eq:ma-coopik-layered}
\end{equation}
两种写法在无奇异、无冲突时结果接近，但在约束冲突时行为不同（\cref{tab:ma-ik-choice}）。

\begin{table}[H]\centering\small
\caption{一次式 vs 分层式协调逆运动学}
\label{tab:ma-ik-choice}
\begin{tabular}{p{2.7cm}p{3.2cm}p{4.0cm}p{3.0cm}}
\toprule
写法 & 绝对/相对关系 & 冲突时行为 & 适用 \\
\midrule
一次式 $\mathbf{J}_{\text{coop}}^{+}[\boldsymbol{\mathcal{V}}_{\text{abs}};\boldsymbol{\mathcal{V}}_{\text{rel}}]$ & 平权（同一伪逆里加权折中） & 绝对与相对误差按伪逆“平摊”，可能都不精确 & 两通道同等重要、追求整体最小二乘 \\
分层式 $\mathbf{P}_{\text{rel}}\mathbf{J}_{\text{abs}}^{+}\boldsymbol{\mathcal{V}}_{\text{abs}}$ & 相对优先 & 相对约束\emph{严格优先}满足，绝对运动让步 & 内力安全 $>$ 物体精度（刚体把持默认） \\
\bottomrule
\end{tabular}
\end{table}

\begin{insight}[刚体把持应优先用分层式（相对约束优先）]
“不挤压物体”（相对约束）是\emph{硬安全}约束，“物体精确到位”（绝对）是\emph{可让步}的性能目标。一次式把两者平权折中，冲突时可能轻微违反相对约束（产生内力）以换取绝对精度，这在易碎物上不可接受。这与机械臂 WBC 的任务优先级（分层 QP）是同一思想：安全/约束类任务放高优先级，性能类任务放低优先级，低优先级在高优先级零空间里执行。$N$ 臂时这进一步泛化为 \cref{sec:ma-allocation} 的“硬约束（$\mathbf{G}\mathbf{h}=\mathbf{F}$、摩擦锥）$+$ 软目标（最小化力）”分层 QP。
\end{insight}

\subsection{非对称任务：扩展相对雅可比}
\label{sec:ma-extjac}

对称分解（绝对 $=$ 中点、相对 $=$ 右减左）隐含一个假设：\emph{两臂地位对等}。但很多真实任务是\emph{非对称}的——一只手“扶”（提供稳定参考），另一只手“操作”（执行精细动作），如左手按食材右手切、左手持工件右手拧螺丝。这类任务里我们关心的不是“物体中点怎么动”，而是操作手相对扶持手的\emph{相对运动}（切/拧）与扶持手自身的\emph{绝对运动}（食材/工件放哪）。Almeida、Karayiannidis 等（2019）提出\textbf{扩展相对雅可比}\cite{almeida2019asymmetric}\pz{存疑: 第二作者应为 Karayiannidis 非"Lima", 会议应为 ISRR 2019 非 IROS; arXiv:1905.01248 与内容归属确, 见 \cref{tab:ma-history}}：把“相对运动”和“指定的某一只手的绝对运动”组合进一个雅可比，
\begin{equation}
  \begin{bmatrix}\boldsymbol{\mathcal{V}}_{\text{rel}}\\ \boldsymbol{\mathcal{V}}_{\text{abs}}^{(k)}\end{bmatrix}=\mathbf{J}_{\text{ext}}\,\dot{\mathbf{q}},
  \qquad
  \mathbf{J}_{\text{ext}}=\begin{bmatrix}\mathbf{J}_{\text{rel}}\\ \mathbf{J}_{\text{abs}}^{(k)}\end{bmatrix},
  \label{eq:ma-extjac}
\end{equation}
其中 $\boldsymbol{\mathcal{V}}_{\text{abs}}^{(k)}$ 不再是两手中点，而是\emph{所选参考手 $k$}（如扶持手）的绝对 twist。这相当于把对称表述里“绝对 $=$ 中点”替换成“绝对 $=$ 某只指定手”——一个\emph{可调的参考权重}：
\begin{equation}
  \boldsymbol{\mathcal{V}}_{\text{abs}}^{(\alpha)}=\alpha\,\bar{\boldsymbol{\mathcal{V}}}_L+(1-\alpha)\,\bar{\boldsymbol{\mathcal{V}}}_R.
  \label{eq:ma-alpha}
\end{equation}
$\alpha=\tfrac12$ 退化为 Uchiyama--Dauchez 对称中点；$\alpha=1$ 表示“以左手为绝对参考”（左手扶、右手操作）；$\alpha=0$ 反之。这个 $\alpha$ 把对称/非对称\emph{统一进一个连续参数}（\cref{sec:ma-symmetry} 详述其选取）。

\begin{insight}[从对称到非对称是“集体运动以谁为代表”的选择]
从对称（$\alpha=\tfrac12$）到非对称（$\alpha\in\{0,1\}$）的连续过渡，本质是“绝对运动这个集体坐标，以谁为代表”的选择。对称表述把集体坐标定义为民主的“平均”，非对称表述把它定义为“指定代表（leader）”——这在 \cref{sec:ma-arch} 会直接对应“对称协调 vs 主从架构”的体系结构选择。运动学层面的 $\alpha$ 选择，预示了控制架构层面的主从 vs 对称之争。
\end{insight}

\begin{table}[H]\centering\small
\caption{协同运动学：单系统视角 vs 本章多机视角的分工}
\label{tab:ma-kin-division}
\begin{tabular}{p{3.0cm}p{5.4cm}p{5.0cm}}
\toprule
主题 & 单系统视角（机械臂方向） & 本章多机视角 \\
\midrule
增广雅可比 $\mathbf{J}_{\text{aug}}$ & 详细推导 $+$ Pinocchio C++/Python 实现 & 仅作符号，引用之 \\
相对雅可比 $\mathbf{J}_{\text{rel}}$ & 给出 $\bar{\mathbf{J}}_R-\bar{\mathbf{J}}_L$ 工程形式 $+$ 平移算子细节 & 从对称表述第一性原理重新推导其来历 \\
绝对雅可比 $\mathbf{J}_{\text{abs}}$ & 未强调 & 作为内力控制的对偶搭档引入 \\
统一参数 $\alpha$ & 未涉及 & 引入，连接 \cref{sec:ma-symmetry} 架构选择 \\
零空间三大用途 & 详细（限位/操作度/自碰撞回避） & 引用，不重复 \\
Jamisola--Roberts 紧凑形式 & 提及 & 给出工程拼装结论 \\
\bottomrule
\end{tabular}
\end{table}

\begin{pitfall}[编程陷阱：用 \texttt{LOCAL\_WORLD\_ALIGNED} 雅可比直接做 $\mathbf{J}_R-\mathbf{J}_L$ 当相对雅可比]
现象：物体绕自身质心整体转动（合法的绝对运动）时，算出的 $\boldsymbol{\mathcal{V}}_{\text{rel}}\neq\mathbf{0}$，控制器误以为有相对运动 $\to$ 误施加内力 $\to$ 物体被无故挤压。根本原因：该约定只对齐了坐标轴方向，没把线速度参考点移到同一点；刚体转动时两末端线速度天然不同，不统一参考点就做差，会把“绕物体转”误算成“相对平移”。正确做法：先用平移算子 $\mathbf{X}(\mathbf{p}_c-\mathbf{p}_i)$ 把两臂 twist 统一到公共参考点，再做 $\mathbf{J}_{\text{rel}}=\bar{\mathbf{J}}_R-\bar{\mathbf{J}}_L$。自检：构造让物体绕质心纯转动的 $\dot{\mathbf{q}}$，验证 $\|\mathbf{J}_{\text{rel}}\dot{\mathbf{q}}\|\approx0$；若显著非零，说明参考点没统一对。
\end{pitfall}

\begin{pitfall}[概念误区：认为“绝对雅可比 $\mathbf{J}_{\text{abs}}$ 没用，有 $\mathbf{J}_{\text{rel}}$ 投影就够”]
新手想法：“刚体把持只要把运动投影到 $\ker(\mathbf{J}_{\text{rel}})$ 保证不产生内力就行，不需要单独的 $\mathbf{J}_{\text{abs}}$。”实际上：投影到 $\ker(\mathbf{J}_{\text{rel}})$ 只保证“不破坏相对位姿”，但你还得在这个零空间里实现期望的物体运动——而“物体怎么动”恰恰是 $\mathbf{J}_{\text{abs}}$ 描述的。没有 $\mathbf{J}_{\text{abs}}$，你能保证不挤压物体，却无法精确控制物体走到哪；更重要的是，内力控制（\cref{sec:ma-internalforce}）需要 $\mathbf{J}_{\text{abs}}/\mathbf{J}_{\text{rel}}$ 的对偶（力分配），缺了 $\mathbf{J}_{\text{abs}}$ 内/外力解耦无从谈起。二者是对称表述的一对孪生坐标，缺一不可。
\end{pitfall}

\begin{exercise}
\begin{enumerate}
  \item \textbf{[推导]} 从相对位姿 $\mathbf{T}_{\text{rel}}=\mathbf{T}_L^{-1}\mathbf{T}_R$ 出发对时间求导（用 $\SE(3)$ 上的左/右平凡化），导出相对 twist $\boldsymbol{\mathcal{V}}_{\text{rel}}$ 关于 $\boldsymbol{\mathcal{V}}_L,\boldsymbol{\mathcal{V}}_R$ 的表达式，并说明伴随变换与本节平移算子 $\mathbf{X}(\cdot)$ 在什么近似下一致（提示：当两末端姿态相同、仅位置不同时，伴随退化为平移算子）。
  \item \textbf{[验证]} 为双 $7$ 自由度臂计算 $\mathbf{J}_{\text{abs}},\mathbf{J}_{\text{rel}}$，构造三种关节速度验证语义：(a) 物体纯平移（$\boldsymbol{\mathcal{V}}_{\text{rel}}=\mathbf{0}$）；(b) 物体绕质心纯转动（$\boldsymbol{\mathcal{V}}_{\text{rel}}=\mathbf{0}$）；(c) 两手相对拉伸（$\boldsymbol{\mathcal{V}}_{\text{rel}}\neq\mathbf{0}$，对应内力）。记录每种 $\|\boldsymbol{\mathcal{V}}_{\text{abs}}\|,\|\boldsymbol{\mathcal{V}}_{\text{rel}}\|$。
  \item \textbf{[设计]} 把 $\alpha$ 从 $0$ 连续扫到 $1$，对“左手持杯、右手持壶倒水”任务，分析“以谁为绝对参考”如何影响控制语义。$\alpha$ 取多少最合理？为什么倒水任务里“杯子（左手）”更适合当绝对参考？
  \item \textbf{[跨章综合]} 论证：绝对/相对分解是不是一种“坐标变换式解耦”，和 ADMM 的“变量分裂 $+$ 一致性约束”在思想上有何异同（提示：ADMM 分裂各 agent 的局部变量，绝对/相对分裂任务空间的语义通道）？
\end{enumerate}
\end{exercise}

% ════════════════════════════════════════════════════════════════
\section{闭链约束的两个面孔：$\ker(\mathbf{J}_{\text{rel}})$ 与 $\ker(\mathbf{G})$ 的对偶}
\label{sec:ma-duality}

\paragraph{动机：两套语言，一个物理。}
协同搬运章用 grasp matrix $\mathbf{G}$ 在\emph{力空间}描述协同搬运（内力 $=\ker(\mathbf{G})$）；\cref{sec:ma-coopkin} 用相对雅可比 $\mathbf{J}_{\text{rel}}$ 在\emph{速度空间}描述双臂协调（同步运动 $=\ker(\mathbf{J}_{\text{rel}})$）。这两套语言看似各说各话，但它们描述的是\emph{同一个物理约束}（刚体把持）的两个投影。回顾两套描述。\textbf{力空间}：$N$ 只手对物体施加的 wrench 堆叠为 $\mathbf{h}=[\mathbf{h}_1;\dots;\mathbf{h}_N]$，
\begin{equation}
  \mathbf{F}_{\text{load}}=\mathbf{G}\,\mathbf{h},\qquad
  \mathbf{G}=\begin{bmatrix}\mathbf{G}_1 & \cdots & \mathbf{G}_N\end{bmatrix},\quad
  \mathbf{G}_i=\begin{bmatrix}\mathbf{I}_3 & \mathbf{0}\\ [\mathbf{r}_i]_\times & \mathbf{I}_3\end{bmatrix}\in\mathbb{R}^{6\times6},
  \label{eq:ma-grasp}
\end{equation}
每只手传递完整 $6$ 维 wrench（力 $+$ 力矩，因夹爪刚性把持能传力矩；协同搬运章的点接触版本 $\mathbf{G}_i\in\mathbb{R}^{6\times3}$ 只传力，是这里的退化），$\mathbf{r}_i$ 是第 $i$ 抓取点到物体质心的向量，内力 $=\ker(\mathbf{G})$、维度 $6N-6$。\textbf{速度空间}：刚体把持要求两手相对位姿不变，即 $\boldsymbol{\mathcal{V}}_{\text{rel}}=\mathbf{J}_{\text{rel}}\dot{\mathbf{q}}=\mathbf{0}$，满足约束的关节速度 $=\ker(\mathbf{J}_{\text{rel}})$。问题：这两个零空间是什么关系？它们都和“内力”有关，但一个在力里、一个在速度里。

\subsection{运动学-静力学对偶：虚功原理}
\label{sec:ma-vwork}

答案藏在\textbf{虚功原理}里。考虑物体侧：物体运动 twist 是 $\boldsymbol{\mathcal{V}}_{\text{obj}}\in\mathbb{R}^6$，受净 wrench $\mathbf{F}_{\text{load}}$；各手施加 wrench $\mathbf{h}_i$、各手末端运动 twist $\bar{\boldsymbol{\mathcal{V}}}_i$（已统一到物体质心参考点）。各手施加 wrench $\mathbf{h}$、各手末端有虚 twist $\delta\bar{\boldsymbol{\mathcal{V}}}_{\text{c}}$，物体受净 wrench $\mathbf{F}_{\text{load}}$、有虚 twist $\delta\boldsymbol{\mathcal{V}}_{\text{obj}}$。理想约束的功率平衡（手对物体做的功 $=$ 物体获得的功）：
\begin{equation}
  \mathbf{h}^\top\delta\bar{\boldsymbol{\mathcal{V}}}_{\text{c}}=\mathbf{F}_{\text{load}}^\top\delta\boldsymbol{\mathcal{V}}_{\text{obj}}.
  \label{eq:ma-vwork-balance}
\end{equation}
代入 $\mathbf{F}_{\text{load}}=\mathbf{G}\mathbf{h}$ 得 $\mathbf{h}^\top\delta\bar{\boldsymbol{\mathcal{V}}}_{\text{c}}=(\mathbf{G}\mathbf{h})^\top\delta\boldsymbol{\mathcal{V}}_{\text{obj}}=\mathbf{h}^\top\mathbf{G}^\top\delta\boldsymbol{\mathcal{V}}_{\text{obj}}$。对任意 $\mathbf{h}$ 成立，故
\begin{equation}
  \boxed{\;\delta\bar{\boldsymbol{\mathcal{V}}}_{\text{c}}=\mathbf{G}^\top\delta\boldsymbol{\mathcal{V}}_{\text{obj}}\;}
  \label{eq:ma-Gtransp}
\end{equation}
即\emph{接触点速度由物体速度经 $\mathbf{G}^\top$ 得到}。力的映射 $\mathbf{G}$ 和速度的映射 $\mathbf{G}^\top$ 是同一个 $\mathbf{G}$ 的转置——这是运动学（速度）与静力学（力）的对偶。

\begin{insight}[grasp matrix 同时编码运动学与静力学]
$\mathbf{G}$ 不只是“力的映射矩阵”。它同时编码抓取的\emph{运动学}与\emph{静力学}：$\mathbf{G}$（前向）把接触力合成为物体 wrench；$\mathbf{G}^\top$（转置）把物体速度分配为接触点速度。一个矩阵，两个方向，分别管力和速度——这就是“对偶”的精确含义。协同搬运章只用了 $\mathbf{G}$ 的力方向，本章揭示它的速度方向，从而和相对雅可比接上。
\end{insight}

\subsection{两个零空间的物理对应}
\label{sec:ma-nullspaces}

\rebuilt{本小节"两零空间对偶"的核心论断有误, 整段待修(勿原位重建): 本章 grasp 矩阵取 $\mathbf{G}\in\mathbb{R}^{6\times6N}$ 满行秩 6, 故 $\dim\ker(\mathbf{G})=6N-6$(内力, 正确), 但 $\dim\ker(\mathbf{G}^\top)=6-6=0$(一般为平凡), \emph{并非} $6N-6$。所谓"$\ker(\mathbf{G}^\top)$=破坏约束的相对运动、与 $\ker(\mathbf{G})$ 同维同构、由 SVD 配对"系误述: 内力 $\ker(\mathbf{G})$ 的真正对偶是\emph{可控物体运动} $\operatorname{Im}(\mathbf{G}^\top)$(6 维), 破坏约束的相对运动是接触-twist 空间中 $\operatorname{Im}(\mathbf{G}^\top)$ 的正交补(即 $\ker$ 于速度映射, 维 $6N-6$), 而非 $\ker(\mathbf{G}^\top)$; 非平凡的 $\ker(\mathbf{G}^\top)$ 反而是抓取"不可定"的退化情形。可救结论: $\dim\ker(\mathbf{G})=6N-6$ 与 \cref{eq:ma-final-equiv} 的双条件($\mathbf{J}_{\text{rel}}\dot{\mathbf{q}}=\mathbf{0}\iff\bar{\boldsymbol{\mathcal{V}}}_{\text{c}}\in\operatorname{Im}(\mathbf{G}^\top)$)成立; 其余"$\ker(\mathbf{G}^\top)$ 对偶"叙述均待重写。}%
现在可精确回答“$\ker(\mathbf{G})$ 和 $\ker(\mathbf{J}_{\text{rel}})$ 是什么关系”。\textbf{$\ker(\mathbf{G})$（力的零空间，维度 $6N-6$）}：满足 $\mathbf{G}\mathbf{h}=\mathbf{0}$ 的 wrench 组合——它们互相抵消、不推动物体，只在物体内部“较劲”，即内力。\textbf{$\operatorname{Im}(\mathbf{G}^\top)$（速度的值空间，维度 $6$）}：所有“由某物体运动产生的接触点速度模式”——协调一致的运动，各手速度恰对应同一刚体运动。\textbf{$\ker(\mathbf{G}^\top)$}：\pz{待修(原断言已删): 源称其维度 $6N-6$、代表"破坏约束的相对运动"——均误; $\ker(\mathbf{G}^\top)$ 一般为平凡(维 0), 破坏约束的相对运动应为 $\operatorname{Im}(\mathbf{G}^\top)$ 的正交补, 详见本小节开头 \texttt{重建待核对} 说明}。对双臂（$N=2$）整理为 \cref{tab:ma-subspaces}（其中 $\ker(\mathbf{G}^\top)$ 行待修）。

\begin{table}[H]\centering\small
\caption{双臂（$N=2$）闭链子空间的物理含义}
\label{tab:ma-subspaces}
\begin{tabular}{llll}
\toprule
子空间 & 维度（$N=2$） & 物理含义 & 在速度还是力 \\
\midrule
$\ker(\mathbf{G})$ & $12-6=6$ & 内力（两手互相挤压/拉扯，不动物体） & 力 \\
$\operatorname{Im}(\mathbf{G}^\top)$ & $6$ & 协调运动（两手速度对应同一刚体运动） & 速度 \\
% 【待修·原行有误，已注释】下行把 ker(G^T) 标为 6 维"相对运动"系误述：
% $\ker(\mathbf{G}^\top)$ & $12-6=6$ & 相对运动（破坏刚性把持的两手相对速度） & 速度 \\
$\ker(\mathbf{G}^\top)$ & $6-6=0$（一般平凡） & \rebuilt{待修: 原标 6 维"相对运动"误; 破坏约束的相对运动是 $\operatorname{Im}(\mathbf{G}^\top)$ 在接触-twist 空间的正交补(维 6), 非 $\ker(\mathbf{G}^\top)$} & 速度 \\
\bottomrule
\end{tabular}
\end{table}

\noindent\textbf{关键对应}\rebuilt{待修(勿原位重建): 下述"$\ker(\mathbf{G}^\top)$ 与 $\ker(\mathbf{G})$ 同维同构、由 SVD 配对"系误述, 对应公式已注释。正确方向: 内力 $\ker(\mathbf{G})$(维 $6N-6$)与"破坏约束的相对运动"(接触-twist 空间中 $\operatorname{Im}(\mathbf{G}^\top)$ 的正交补, 同维 $6N-6$)配对; 该相对运动空间不是 $\ker(\mathbf{G}^\top)$(后者一般平凡), 配对也不来自 $\mathbf{G}$ 的 SVD。}：%
% 【待修·原断言有误，整段注释】源称 ker(G^T)(6维)与 ker(G)(6维)同维同构、经 SVD 配对——误：
% 速度里“破坏约束的相对运动方向”（$\ker(\mathbf{G}^\top)$，$6$ 维）与力里“内力方向”（$\ker(\mathbf{G})$，$6$ 维）\emph{同维数}，且通过 $\mathbf{G}$ 的奇异值分解自然配对——它们是同构的。这正是为什么内力总和相对运动配对：
% \begin{equation}
%   \underbrace{\text{相对运动}}_{\ker(\mathbf{G}^\top),\ 6\text{维}}\;\xleftrightarrow{\ \text{对偶}\ }\;\underbrace{\text{内力}}_{\ker(\mathbf{G}),\ 6\text{维}}.
%   \label{eq:ma-dualpair}
% \end{equation}
\noindent 内力与"破坏约束的相对运动"在双臂下均为 $6$ 维、互为对偶（具体配对的正确叙述见上 \texttt{重建待核对}），这为 \cref{sec:ma-internalforce} 用相对位姿偏差闭环控制内力提供依据。

\begin{insight}[内力和相对运动是同一枚硬币的两面]
物理直觉：如果你能让两手产生某个“破坏刚性约束的相对运动”（拉伸），那么阻止这个运动（刚体不让你拉伸）就会产生对应方向的“内力”。\textbf{能拉伸的方向 $=$ 能产生内力的方向}，由抓取几何 $\mathbf{G}$ 的零空间结构锁定。这就是为什么 \cref{sec:ma-internalforce} 控制内力时，控制的变量是“相对位姿偏差”——因为相对位姿偏差正是内力的“位移共轭”。
\end{insight}

\subsection{把 $\mathbf{J}_{\text{rel}}$ 接进来：关节空间的完整图景}
\label{sec:ma-jointpicture}

上面是“物体侧”的对偶（接触点 $\leftrightarrow$ 物体）。还差一步：把它连到\emph{关节空间}（我们真正能控制的）。关节速度 $\dot{\mathbf{q}}$ 通过各臂雅可比产生末端速度 $\bar{\boldsymbol{\mathcal{V}}}_{\text{c}}=\mathbf{J}_{\text{aug}}\dot{\mathbf{q}}$（增广雅可比）；末端速度要满足“对应某个刚体运动”才不破坏抓取。把它分解到“协调运动子空间 $\operatorname{Im}(\mathbf{G}^\top)$”和“破坏约束子空间（$\operatorname{Im}(\mathbf{G}^\top)$ 的正交补）”\pz{待修: 源此处把该正交补记作 $\ker(\mathbf{G}^\top)$ 系误(后者一般平凡); 正确为 $\operatorname{Im}(\mathbf{G}^\top)$ 在接触-twist 空间的正交补, 维 $6N-6$, 详见 \cref{sec:ma-nullspaces} 开头说明}：落在 $\operatorname{Im}(\mathbf{G}^\top)$ 的分量对应物体运动 $\boldsymbol{\mathcal{V}}_{\text{obj}}$（绝对运动，$\boldsymbol{\mathcal{V}}_{\text{abs}}$ 是它的一个 $6$ 维参数化），落在该正交补的分量正是相对运动 $\boldsymbol{\mathcal{V}}_{\text{rel}}=\mathbf{J}_{\text{rel}}\dot{\mathbf{q}}$。于是
\begin{equation}
  \boxed{\;\mathbf{J}_{\text{rel}}\dot{\mathbf{q}}=\mathbf{0}
  \iff \bar{\boldsymbol{\mathcal{V}}}_{\text{c}}\in\operatorname{Im}(\mathbf{G}^\top)
  \iff \text{末端速度对应某个纯刚体运动，不产生相对运动}\;}
  \label{eq:ma-final-equiv}
\end{equation}
这就是 $\ker(\mathbf{J}_{\text{rel}})$ 与 $\ker(\mathbf{G})$ 对偶的最终形式：$\ker(\mathbf{J}_{\text{rel}})$（关节速度满足刚体约束）映射到末端的 $\operatorname{Im}(\mathbf{G}^\top)$（协调运动），而它的“互补”——会产生 $\boldsymbol{\mathcal{V}}_{\text{rel}}\neq\mathbf{0}$ 的关节速度——映射到\emph{破坏约束的相对运动子空间}（$\operatorname{Im}(\mathbf{G}^\top)$ 的正交补，维 $6N-6$）\pz{待修: 源此处记作 $\ker(\mathbf{G}^\top)$ 系误, 见 \cref{sec:ma-nullspaces} 开头}，后者与 $\ker(\mathbf{G})$（内力）对偶。一句话串起前三节：

\begin{insight}[一句话串起前三节]
物体运动（绝对，$\boldsymbol{\mathcal{V}}_{\text{abs}}$，$\operatorname{Im}(\mathbf{G}^\top)$）和内力（$\ker(\mathbf{G})$）是两个互补的、解耦的控制通道；相对运动（$\boldsymbol{\mathcal{V}}_{\text{rel}}$，落在 $\operatorname{Im}(\mathbf{G}^\top)$ 的正交补\pz{待修: 源记作 $\ker(\mathbf{G}^\top)$ 系误, 见 \cref{sec:ma-nullspaces} 开头}）是内力的速度共轭，控制内力就是控制相对位姿偏差。
\end{insight}

\begin{insight}[$\mathbf{J}_{\text{rel}}$ 与 $\mathbf{G}$ 的对偶 vs 电路的基尔霍夫定律（带边界）]
\textbf{像的部分}：KCL（节点电流和为零）约束“流量”（类比速度/运动），KVL（回路电压和为零）约束“势差”（类比力/内力），二者描述同一电路拓扑的两个对偶方面；闭链的速度约束 $\mathbf{J}_{\text{rel}}\dot{\mathbf{q}}=\mathbf{0}$ 和力约束 $\mathbf{G}\mathbf{h}=\mathbf{F}_{\text{load}}$ 也描述同一机械回路的两个对偶方面，由同一几何（抓取拓扑）决定。\textbf{不像的部分}：电路对偶线性恒定；机械臂的 $\mathbf{G}$ 和 $\mathbf{J}_{\text{rel}}$ 随构型变化（$\mathbf{r}_i$、雅可比都依赖 $\mathbf{q}$），对偶关系在每个瞬时重新建立。\textbf{不要延伸到}：“电路里电流电压完全独立可测，所以内力和运动也完全独立可控”——机械系统里相对运动一旦发生就立即通过物体刚度产生内力，二者动力学上紧密耦合，“解耦”是控制器每个周期主动维持的，不是天然独立。
\end{insight}

\subsection{最简双臂把对偶算出来}
\label{sec:ma-worked-duality}

用一个可手算的最简实例把对偶“落地”。设两手沿 $x$ 轴对称把持一根杆，物体质心在原点，左手在 $\mathbf{r}_L=(-\ell,0,0)$、右手 $\mathbf{r}_R=(+\ell,0,0)$，$\ell=0.15$\,m。只考虑平面内的力（不含力矩，用 $3$ 维 wrench 简化，结论同样成立）：
\begin{equation}
  \mathbf{G}=\begin{bmatrix}\mathbf{G}_L & \mathbf{G}_R\end{bmatrix},\quad
  \mathbf{G}_i=\begin{bmatrix}1&0\\0&1\\-r_{iy}&r_{ix}\end{bmatrix}
  \Rightarrow
  \mathbf{G}=\begin{bmatrix}1&0&1&0\\0&1&0&1\\0&-\ell&0&\ell\end{bmatrix}\in\mathbb{R}^{3\times4},
  \label{eq:ma-worked-G}
\end{equation}
其中 $\mathbf{h}_i=(f_{ix},f_{iy})$、物体 wrench $\mathbf{F}=(F_x,F_y,M_z)$，第三行是力矩 $\mathbf{r}\times\mathbf{f}$ 的 $z$ 分量。$\mathbf{G}$ 秩 $3$，故 $\dim\ker(\mathbf{G})=4-3=1$——一维内力。\textbf{求内力方向}：解 $\mathbf{G}\mathbf{h}=\mathbf{0}$。第一行 $f_{Lx}+f_{Rx}=0$；第三行 $-\ell f_{Ly}+\ell f_{Ry}=0$ 给 $f_{Ly}=f_{Ry}$，配合第二行 $f_{Ly}+f_{Ry}=0$ 得 $f_{Ly}=f_{Ry}=0$。于是 $\ker(\mathbf{G})=\operatorname{span}\{(1,0,-1,0)\}$——即 $f_{Lx}=+1,f_{Rx}=-1$：\emph{左手向右推、右手向左推，两手沿杆轴相向对夹}。这正是物理直觉里的“夹紧内力”，现在它是 $\ker(\mathbf{G})$ 的一个具体基向量。\textbf{验证它不动物体}：$\mathbf{G}\cdot(1,0,-1,0)^\top=(1-1,0,0)^\top=\mathbf{0}$。\textbf{对偶的相对运动方向}：内力沿 $x$ 轴对夹 $\leftrightarrow$ 速度上“破坏约束的相对运动”也沿 $x$ 轴——两手沿杆轴相对拉伸/压缩。一维内力 $\leftrightarrow$ 一维相对拉伸，同维同构，对偶在这个最简例子上\emph{显式成立}。

\begin{insight}[从实例回看抽象]
这个 $3\times4$ 的小 $\mathbf{G}$ 把整节抽象对偶压缩成“$(1,0,-1,0)$ 这个向量”——它既是 $\ker(\mathbf{G})$ 的基（力空间的内力），又指向相对拉伸方向（速度空间）。当你在 $N$ 臂、$6$D wrench 的大矩阵里迷失时，回到这个例子：内力永远是“某种相向的力组合，合成为零、纯内部较劲”，相对运动永远是“破坏刚性约束的那部分末端运动”，二者由 $\mathbf{G}$ 的零空间结构一一配对。维度变大只是基变多，结构不变。
\end{insight}

\subsection{从对偶看“为什么内力必须闭环”}
\label{sec:ma-why-closedloop}

对偶视角直接解释一个工程上反直觉的事实：\textbf{内力不能开环施加，必须闭环维护}（这是 \cref{sec:ma-internalforce} 的核心，这里先用对偶把动机讲透）。设想开环：你想要内力 $\mathbf{f}_{\text{int}}$，在 $\ker(\mathbf{G})$ 里选一个 wrench，让两手按它出力。但 $\ker(\mathbf{G})$ 依赖 $\mathbf{r}_i$（抓取点到质心向量），$\mathbf{r}_i$ 依赖物体当前位姿；物体一动 $\to\mathbf{r}_i$ 变 $\to\ker(\mathbf{G})$ 的基旋转 $\to$ 你“开环”算好的 $\mathbf{f}_{\text{int}}$ 不再落在当前 $\ker(\mathbf{G})$ 里 $\to$ 它有了沿 $\operatorname{Im}$ 方向的分量 $\to$ 这个分量推动物体！本想“只夹紧不推物体”的内力，因为几何变了，变成“既夹紧又乱推”，物体轨迹被污染、误差累积、失控。正确做法（\cref{sec:ma-internalforce}）：闭环测量相对位姿偏差（内力的位移共轭），用阻抗律 $\mathbf{f}_{\text{int}}=\mathbf{K}_{\text{int}}\cdot(\text{相对位姿偏差})$ 实时产生内力，几何变了控制器自动跟上。

\begin{insight}[内力必须闭环，因为 $\ker(\mathbf{G})$ 本身随构型旋转]
开环把力扔进一个会动的子空间，力会“漏”到运动通道里。闭环（用相对位姿偏差反馈）相当于始终把力对准当前的 $\ker(\mathbf{G})$，无论它怎么转。这与共识“需要持续通信而非一次性广播”异曲同工——环境/拓扑在变，必须持续校正而非开环下发。
\end{insight}

\begin{pitfall}[概念误区：认为“内力是浪费，应该尽量调成零”]
新手想法：“内力不推动物体，纯属内耗，最优控制应该让内力 $=0$。”实际上：内力是抓取稳定性的来源。对悬空刚体，若内力为零，两手只靠摩擦维持抓取——稍有扰动（物体惯性、外部碰撞）就打滑脱落。正的内力（两手相向“夹紧”）提供摩擦裕度，是必需的。内力的目标不是零，而是“足够大以防滑、足够小以不压坏物体”的一个区间（\cref{sec:ma-internalforce} 给设计准则）。
\end{pitfall}

\begin{pitfall}[思维陷阱：把 $\ker(\mathbf{G})$ 当成固定子空间]
新手想法：“grasp matrix 的零空间是几何决定的常数子空间。”实际上：$\mathbf{G}$ 依赖各接触点到质心的向量 $\mathbf{r}_i$，物体一旋转/平移 $\mathbf{r}_i$ 改变、$\mathbf{G}$ 改变、$\ker(\mathbf{G})$ 的基随之旋转。把 $\ker(\mathbf{G})$ 当常数会导致开环内力泄漏到运动通道。正确思维：$\ker(\mathbf{G})$ 是随构型变化的“活”子空间，每个控制周期重算。
\end{pitfall}

\begin{exercise}
\begin{enumerate}
  \item \textbf{[推导]} 用虚功原理补全 $\delta\bar{\boldsymbol{\mathcal{V}}}_{\text{c}}=\mathbf{G}^\top\delta\boldsymbol{\mathcal{V}}_{\text{obj}}$ 每一步的理由，并说明：为什么“理想约束虚功为零”对刚性把持成立、对有摩擦滑动的接触不成立？
  \item \textbf{[计算]} 对双臂（$N=2$）刚性把持分别计算 $\dim\ker(\mathbf{G})$、$\dim\operatorname{Im}(\mathbf{G}^\top)$、$\dim\ker(\mathbf{G}^\top)$，验证维度定理；再对三臂（$N=3$）重做，观察内力维度如何随 $N$ 增长。
  \item \textbf{[编程]} 为双臂共持物体的某构型同时计算 $\mathbf{G}$（从抓取点几何）和 $\mathbf{J}_{\text{rel}}$（从雅可比），数值验证：物体绕质心纯转动的 $\dot{\mathbf{q}}$ 满足 $\mathbf{J}_{\text{rel}}\dot{\mathbf{q}}\approx\mathbf{0}$，且对应末端速度 $\mathbf{J}_{\text{aug}}\dot{\mathbf{q}}$ 落在 $\operatorname{Im}(\mathbf{G}^\top)$ 中。
  \item \textbf{[思考]} $\boldsymbol{\mathcal{V}}_{\text{rel}}$ 和 $\mathbf{f}_{\text{int}}$ 维度都是 $6$。论证它们构成一对“功率共轭变量”（$P_{\text{int}}=\mathbf{f}_{\text{int}}^\top\boldsymbol{\mathcal{V}}_{\text{rel}}$ 是内力做的功率），并解释为什么刚体把持时 $\boldsymbol{\mathcal{V}}_{\text{rel}}=\mathbf{0}$ 意味着内力不做功（既不输入也不耗散能量）。
\end{enumerate}
\end{exercise}

% ════════════════════════════════════════════════════════════════
\section{内力控制：从双臂到 $N$ 臂}
\label{sec:ma-internalforce}

\paragraph{动机：内力的精确定义与期望值设计。}
\cref{sec:ma-duality} 证明内力（$\ker(\mathbf{G})$）和相对运动（$\ker(\mathbf{G}^\top)$）对偶、且内力必须闭环。本节给出具体控制律。对双臂（$N=2$），各手 wrench $\mathbf{h}_L,\mathbf{h}_R\in\mathbb{R}^6$（已表达在物体质心参考点），物体净 wrench（运动力，质心参考点下 $\mathbf{G}=[\mathbf{I}_6\ \mathbf{I}_6]$）为 $\mathbf{F}_{\text{load}}=\mathbf{G}\mathbf{h}=\mathbf{h}_L+\mathbf{h}_R$。内力定义为两手 wrench 的“反对称部分”——它合成到物体为零、纯内耗：
\begin{equation}
  \boxed{\;\mathbf{f}_{\text{int}}=\tfrac12\big(\mathbf{h}_R-\mathbf{h}_L\big)\;}
  \label{eq:ma-fint}
\end{equation}
验证它落在 $\ker(\mathbf{G})$：把 $\mathbf{h}_L=\tfrac12\mathbf{F}_{\text{load}}-\mathbf{f}_{\text{int}}$、$\mathbf{h}_R=\tfrac12\mathbf{F}_{\text{load}}+\mathbf{f}_{\text{int}}$ 代回 $\mathbf{G}\mathbf{h}=\mathbf{h}_L+\mathbf{h}_R=\mathbf{F}_{\text{load}}$，内力项相消，确实不贡献净 wrench。这给出双臂 wrench 的标准分解：
\begin{equation}
  \mathbf{h}_L=\tfrac12\mathbf{F}_{\text{load}}-\mathbf{f}_{\text{int}},\qquad
  \mathbf{h}_R=\tfrac12\mathbf{F}_{\text{load}}+\mathbf{f}_{\text{int}}.
  \label{eq:ma-hLR}
\end{equation}
物理图景：$\tfrac12\mathbf{F}_{\text{load}}$ 是每只手分担的“搬运负载”（平分），$\pm\mathbf{f}_{\text{int}}$ 是叠加其上的“对夹力”——右手 $+\mathbf{f}_{\text{int}}$（向左推）、左手 $-\mathbf{f}_{\text{int}}$（向右推），两手相向夹紧。\textbf{内力期望值}是个工程权衡，下界由防滑决定、上界由物体强度决定（\cref{tab:ma-fint-bounds}）。

\begin{table}[H]\centering\small
\caption{内力期望值的约束（防滑下界、强度上界）}
\label{tab:ma-fint-bounds}
\begin{tabular}{lll}
\toprule
约束 & 不等式 & 含义 \\
\midrule
防滑下界 & $f_{\text{int}}^{\text{normal}}\ge \|\mathbf{F}_{\text{tangent}}\|/\mu$ & 法向夹紧力须大到让摩擦力足以平衡切向负载（摩擦锥） \\
抗扰裕度 & $f_{\text{int}}\ge f_{\text{slip}}\,(1+\gamma)$ & 留 $\gamma$（如 $50\%$）裕度应对惯性力、外部碰撞 \\
强度上界 & $f_{\text{int}}\le f_{\text{crush}}$ & 不超过物体被压坏的阈值（玻璃、纸盒远小于钢件） \\
电机上界 & $\|\mathbf{h}_i\|\le\tau_{\max,i}$ 经雅可比映射 & 各臂力矩不饱和 \\
\bottomrule
\end{tabular}
\end{table}

\begin{insight}[内力期望值是一个区间而非一个点]
内力期望值不是一个点，而是由“防滑下界”和“强度上界”夹出的\emph{区间}——这是它与运动目标（一个精确轨迹点）的根本不同。运动控制追求“精确到达”，内力控制追求“落在安全区间”。把内力当成需要精确跟踪的硬目标是常见误区；它更像一个带软约束的调节量。这也解释了为什么内力控制天然适合用阻抗（柔顺、容许偏差）而非位置式硬跟踪。
\end{insight}

\textbf{把内力区间算成具体数字。} 双臂水平端持一块 $m=2$\,kg 的玻璃板，两侧面摩擦把持，$\mu=0.4$，被压坏阈值 $f_{\text{crush}}=120$\,N。静止时切向负载 $F_{\text{tangent}}=mg=2\times9.81\approx19.6$\,N（重力沿切向、靠摩擦托住）；搬运加速时（设 $a=3$\,m/s$^2$）可达 $\approx26$\,N。防滑下界 $f_{\text{int}}^{\text{normal}}\ge F_{\text{tangent}}/\mu=26/0.4=65$\,N；留 $\gamma=50\%$ 裕度 $f_{\text{int}}\ge65\times1.5\approx98$\,N；强度上界 $f_{\text{int}}\le120$\,N。于是\emph{内力可行区间 $[98,120]$\,N（很窄！）}。

\begin{insight}[区间可能很窄甚至为空]
上面算出的可行区间只有 $[98,120]$\,N，宽度 $22$\,N——留给控制器的余地很小。如果玻璃更脆（$f_{\text{crush}}=90$\,N）或摩擦更差（$\mu=0.3\Rightarrow$ 下界升到 $26/0.3\times1.5\approx130$\,N），\emph{区间为空}——意味着“既不打滑又不压碎”在当前摩擦/几何下不可能，必须改抓取方式（加大接触面、换高摩擦衬垫、或用形封闭而非力封闭）。这解释了为什么真实双臂操作里“怎么抓”（grasp planning）和“夹多紧”（内力控制）必须联合设计：内力控制器再好，也救不了一个可行区间为空的抓取。这也是 \cref{sec:ma-symmetry} 末“角色/抓取选择”与本节内力控制耦合的根源。
\end{insight}

\subsection{Caccavale 内/外阻抗双层结构}
\label{sec:ma-caccavale}

Caccavale 等（2008）的核心思想\cite{caccavale2008sixdof}：用\emph{两个阻抗}分别管两个解耦通道（正是 \cref{sec:ma-coopkin,sec:ma-duality} 的绝对/相对、运动/内力）。\textbf{外环——object impedance（物体级阻抗，管绝对通道）}：让物体对外部环境接触表现出期望柔顺。物体位姿误差 $\tilde{\mathbf{x}}_{\text{obj}}=\mathbf{x}_{\text{obj}}-\mathbf{x}_{\text{obj}}^{\text{des}}$、外部 wrench $\mathbf{F}_{\text{ext}}$：
\begin{equation}
  \mathbf{M}_o\ddot{\tilde{\mathbf{x}}}_{\text{obj}}+\mathbf{D}_o\dot{\tilde{\mathbf{x}}}_{\text{obj}}+\mathbf{K}_o\tilde{\mathbf{x}}_{\text{obj}}=\mathbf{F}_{\text{ext}},
  \label{eq:ma-objimp}
\end{equation}
物体像一个虚拟“质量-弹簧-阻尼”挂在目标轨迹上，外力推它时它柔顺让步——让双臂搬物体去“插孔/贴合”时不会因刚性硬怼而损坏。\textbf{内环——internal impedance（内力阻抗，管相对通道）}：让内力跟踪期望值、同时对相对位姿偏差柔顺。相对位姿误差 $\tilde{\mathbf{x}}_{\text{rel}}$、内力误差 $\tilde{\mathbf{f}}_{\text{int}}=\mathbf{f}_{\text{int}}-\mathbf{f}_{\text{int}}^{\text{des}}$：
\begin{equation}
  \mathbf{M}_i\ddot{\tilde{\mathbf{x}}}_{\text{rel}}+\mathbf{D}_i\dot{\tilde{\mathbf{x}}}_{\text{rel}}+\mathbf{K}_i\tilde{\mathbf{x}}_{\text{rel}}=-\tilde{\mathbf{f}}_{\text{int}},
  \label{eq:ma-intimp}
\end{equation}
把“两手相对位姿”建模成一根虚拟弹簧（刚度 $\mathbf{K}_i$），内力就是弹簧张力。要增大内力（夹紧），就把相对位姿目标设成“两手更靠近”，弹簧被压缩产生夹力——\emph{内力通过相对位姿偏差产生}，正是 \cref{sec:ma-duality}“相对运动是内力的速度共轭”的控制落地。\textbf{两层为什么能解耦}？因为 object impedance 作用在 $\boldsymbol{\mathcal{V}}_{\text{abs}}$ 通道（$\operatorname{Im}(\mathbf{G}^\top)$），internal impedance 作用在 $\boldsymbol{\mathcal{V}}_{\text{rel}}$ 通道（$\ker(\mathbf{G}^\top)$），而 \cref{sec:ma-duality} 证明这两个通道正交。控制器把总关节力矩拆成两部分，分别投影到这两个正交子空间：
\begin{equation}
  \boldsymbol{\tau}=\mathbf{J}_{\text{abs}}^\top\mathbf{F}_{\text{abs}}^{\text{cmd}}+\mathbf{J}_{\text{rel}}^\top\mathbf{F}_{\text{rel}}^{\text{cmd}}+\boldsymbol{\tau}_{\text{null}},
  \label{eq:ma-tau-decomp}
\end{equation}
其中 $\mathbf{F}_{\text{abs}}^{\text{cmd}}$ 来自外环（object impedance $+$ 重力/惯性前馈），$\mathbf{F}_{\text{rel}}^{\text{cmd}}$ 来自内环（internal impedance，产生内力），$\boldsymbol{\tau}_{\text{null}}$ 是冗余自任务（投影到 $\ker(\mathbf{J}_{\text{coop}})$）。

\begin{insight}[Caccavale 双层结构是对偶理论的“控制器化身”]
绝对/相对的运动学分解（\cref{sec:ma-coopkin}）$+$ $\ker(\mathbf{J}_{\text{rel}})/\ker(\mathbf{G})$ 的对偶（\cref{sec:ma-duality}）$\to$ 外环管绝对、内环管相对，两环正交互不干扰。理解了前两节，这个控制器不是“又一个公式”，而是对偶理论的必然产物。反过来，如果不理解对偶，你会困惑“为什么内力用相对位姿弹簧来产生而不是直接命令力”——因为相对位姿是内力的共轭，闭环它才能对抗几何变化（\cref{sec:ma-why-closedloop}）。
\end{insight}

\begin{note}
\textbf{典型刚度量级（双 Franka 共持，工程参考）.}\;\rebuilt{以下增益数量级转述自源稿示例，非定量标定值，仅示意。}外环物体级阻抗常取平移 $\mathbf{K}_o\sim800$\,N/m、旋转 $\sim40$\,N$\cdot$m/rad，阻尼按临界阻尼配；内环内力阻抗刚度更高（$\mathbf{K}_i\sim1500$\,N/m），因为它要在小相对偏差上产生足够内力。控制律的关节力矩按 \cref{eq:ma-tau-decomp} 由 $\boldsymbol{\tau}=\mathbf{J}_{\text{abs}}^\top\mathbf{F}_{\text{abs}}^{\text{cmd}}+\mathbf{J}_{\text{rel}}^\top\mathbf{F}_{\text{rel}}^{\text{cmd}}+\boldsymbol{\tau}_{\text{null}}$ 叠加，所有 wrench 与雅可比须在同一参考点（物体质心）表达。具体 Pinocchio/MuJoCo 实现属工程细节，从略。
\end{note}

\subsection{推广到 $N$ 臂}
\label{sec:ma-narm}

双臂的 $\mathbf{f}_{\text{int}}=\tfrac12(\mathbf{h}_R-\mathbf{h}_L)$ 是 $N=2$ 的特例。$N$ 臂时内力不再是单个 $6$ 维量，而是 $6N-6$ 维子空间里的元素，需要更一般的参数化。回到 grasp matrix：
\begin{equation}
  \mathbf{h}=\mathbf{G}^{+}\mathbf{F}_{\text{load}}+\mathbf{W}\,\boldsymbol{\lambda},\qquad
  \operatorname{Im}(\mathbf{W})=\ker(\mathbf{G}),\ \boldsymbol{\lambda}\in\mathbb{R}^{6N-6},
  \label{eq:ma-narm-param}
\end{equation}
$\mathbf{W}\in\mathbb{R}^{6N\times(6N-6)}$ 的列张成内力空间，$\boldsymbol{\lambda}$ 是内力参数。$N=2$ 时 $\mathbf{W}$ 的一种取法直接对应 $\pm\mathbf{f}_{\text{int}}$：
\begin{equation}
  \mathbf{W}=\tfrac1{\sqrt2}\begin{bmatrix}-\mathbf{I}_6\\ \mathbf{I}_6\end{bmatrix}\in\mathbb{R}^{12\times6}
  \;\Rightarrow\;
  \mathbf{h}=\tfrac12\mathbf{1}\otimes\mathbf{F}_{\text{load}}+\tfrac1{\sqrt2}\begin{bmatrix}-\mathbf{I}_6\\ \mathbf{I}_6\end{bmatrix}(\sqrt2\,\mathbf{f}_{\text{int}}),
  \label{eq:ma-narm-W2}
\end{equation}
恰好还原 \cref{eq:ma-hLR}。$N$ 臂内力的物理结构（Erhart--Hirche\cite{erhart2017internal} 给出严格框架\pz{存疑: 该文年份应为 2015 非 2017; 且"$\binom{N}{2}$ 组成对挤压基"是一般抓取理论的直觉外推, 非该文显式结论(该文反而批评成对/连线式内力刻画无法推广到力矩), 见 \cref{tab:ma-history} 与下注}）：双臂的 $1$ 组成对挤压（两手相向）在 $N$ 臂里成为 $\binom{N}{2}$ 组潜在成对挤压，独立基数从 $6$ 维升到 $6N-6$ 维。

\begin{note}
\textbf{关键细节（Erhart--Hirche, 2015\pz{存疑: 源标 2017, 应为 2015; 且下文"$N\ge3$"限定与"两两配对基"的归属待商, 详见行内批注}）.}\;\pz{待商: 据原文, "grasp 伪逆未必无内力"对任意 $N$ 均成立(原文用 $N=2$ 即给反例), 非仅 $N\ge3$; "两两配对/连线式内力基"是一般抓取文献的视角, 原文恰批评其无法扩展到力矩, 并未把多臂内力表述为"成对配对基"}$N\ge3$ 时“内力”的定义比双臂微妙——简单地用 grasp 伪逆 $\mathbf{G}^{+}$ 得到的力分配\emph{未必无内力}，因为内力的判定依赖末端运动学（接触点不能独立移动）。严格的“无内力广义逆”需把末端运动学约束一并纳入，不能只看 wrench。这是从双臂直觉外推到多臂时最容易踩的坑：双臂的 $\mathbf{f}_{\text{int}}=\tfrac12(\mathbf{h}_R-\mathbf{h}_L)$ 形式简单，让人以为 $N$ 臂也能“两两配对”理解，但多臂内力是一个整体的 $6N-6$ 维子空间，配对只是它的一组（非唯一）基。
\end{note}

\textbf{$N$ 臂内力控制律}：每只手的期望 wrench 为“分担的运动力 $+$ 内力分量”，运动力由物体阻抗外环给出，内力由各对相对位姿的阻抗内环给出（推广双臂内环）：
\begin{equation}
  \mathbf{h}_i^{\text{des}}=\underbrace{(\mathbf{G}^{+})_i\,\mathbf{F}_{\text{load}}^{\text{cmd}}}_{\text{运动力，外环}}+\underbrace{(\mathbf{W}\boldsymbol{\lambda}^{\text{des}})_i}_{\text{内力，内环}},
  \label{eq:ma-narm-ctrl}
\end{equation}
内环的 $\boldsymbol{\lambda}^{\text{des}}$ 通过 QP 或解析在 $\ker(\mathbf{G})$ 里选取（满足摩擦锥、最小化最大归一化力——正是 \cref{sec:ma-allocation} 的内容）。

\begin{insight}[$N$ 臂照搬双臂“$\tfrac12$ 平分运动力”的代价]
双臂平分（每手 $\tfrac12\mathbf{F}$）在两手对称、能力相同时合理。但 $N$ 臂时各手位置（$\mathbf{r}_i$）、能力（$\tau_{\max,i}$）不同，平分会让最弱/最不利位置的手先饱和（\cref{sec:ma-perspective} 反面教材“四臂抬钢板”）。运动力分配应该用 $\mathbf{G}^{+}$（最小范数）或加权伪逆（按能力加权），而非朴素平分——双臂的 $\tfrac12$ 只是 $\mathbf{G}^{+}$ 在对称情形的巧合值。
\end{insight}

\begin{table}[H]\centering\small
\caption{内力控制：单系统视角 vs 本章多机视角的分工}
\label{tab:ma-force-division}
\begin{tabular}{p{3.4cm}p{5.0cm}p{5.0cm}}
\toprule
主题 & 单系统视角 & 本章多机视角 \\
\midrule
内力分解 $\mathbf{f}_{\text{int}}=\tfrac12(\mathbf{h}_R-\mathbf{h}_L)$ & 给出双臂形式 & 从 $\ker(\mathbf{G})$ 推导 $+$ 推广到 $N$ 臂 $\mathbf{W}\boldsymbol{\lambda}$ \\
object/internal impedance & 工程实现细节 & 强调它是绝对/相对对偶的控制化身 \\
无源性 / 时延 / 波变量 & 详细 & 不涉及（属单系统力控工程） \\
$N$ 臂内力的严格定义 & 未涉及（聚焦双臂） & Erhart--Hirche 框架，多臂内力依赖运动学 \\
与摩擦锥/QP 分配接口 & 简述 & 指向 \cref{sec:ma-allocation} 多臂分配 \\
\bottomrule
\end{tabular}
\end{table}

\begin{pitfall}[编程陷阱：内环和外环的 wrench 在不同参考点/坐标系下叠加]
错误做法：object impedance 算出的 $\mathbf{F}_{\text{abs}}$ 在世界系、internal impedance 的 $\mathbf{F}_{\text{rel}}$ 在某末端系，直接 $\mathbf{J}_{\text{abs}}^\top\mathbf{F}_{\text{abs}}+\mathbf{J}_{\text{rel}}^\top\mathbf{F}_{\text{rel}}$ 叠加。现象：物体被搬运时内力随姿态漂移，看似“控制器有 bug 但找不到”。根本原因：两个通道的 wrench 必须在同一参考点（通常物体质心）、同一坐标系表达后才能正确投影回关节空间，$\mathbf{J}_{\text{abs}}/\mathbf{J}_{\text{rel}}$ 也必须对应同一参考点。正确做法：统一所有 wrench 和雅可比到物体质心参考点，内外环一致。
\end{pitfall}

\begin{pitfall}[概念误区：认为“内力阻抗刚度越大越好（夹得越紧越稳）”]
新手想法：“internal impedance 的 $\mathbf{K}_i$ 调大，内力大，抓取就越稳。”实际上：$\mathbf{K}_i$ 过大有两个问题——(1) 内力随相对位姿偏差线性增长，$\mathbf{K}_i$ 大时微小跟踪误差就产生巨大内力 $\to$ 压坏物体或力矩饱和；(2) 高刚度内环降低系统对外部扰动的柔顺性，与外环 object impedance 的“柔顺”目标矛盾，可能引发两环对抗、振荡。正确理解：$\mathbf{K}_i$ 要匹配物体刚度和期望内力区间——通常通过“期望内力 / 可接受相对偏差”反推。
\end{pitfall}

\begin{pitfall}[思维陷阱：把 $N$ 臂内力当成“$N-1$ 组双臂内力的简单拼接”]
新手想法：“三臂就是（臂 1, 臂 2）$+$（臂 2, 臂 3）两组双臂内力。”实际上：$N$ 臂内力是一个整体的 $6N-6$ 维子空间，“两两配对”只是它的一组基，且配对基之间不独立（臂 2 同时出现在两对里）；更重要的是，多臂内力的严格定义依赖末端运动学（Erhart--Hirche），不能只看 wrench 配对。正确思维：用 grasp 零空间 $\ker(\mathbf{G})$ 的整体参数化 $\mathbf{W}\boldsymbol{\lambda}$ 理解 $N$ 臂内力，配对只是直觉辅助。
\end{pitfall}

\begin{exercise}
\begin{enumerate}
  \item \textbf{[推导]} 从 $\mathbf{h}=\mathbf{G}^{+}\mathbf{F}_{\text{load}}+\mathbf{W}\boldsymbol{\lambda}$ 出发，对双臂（质心参考点 $\mathbf{G}=[\mathbf{I}_6\ \mathbf{I}_6]$）求出 $\mathbf{G}^{+}$ 和一组 $\ker(\mathbf{G})$ 的基 $\mathbf{W}$，验证它还原 \cref{eq:ma-hLR}，并说明 $\mathbf{G}^{+}$ 给出的运动力分配为什么是“平分”。
  \item \textbf{[实现]} 实现双臂协调阻抗（Caccavale 内/外双层），让双臂共持一根弹性梁，跟踪一条物体圆弧轨迹同时维持 $20$\,N 内力，调 $\mathbf{K}_i,\mathbf{D}_i$ 使内力波动 $<10\%$。
  \item \textbf{[设计扩展]} 把双臂协调阻抗推广到三臂：实现 grasp matrix $\mathbf{G}\in\mathbb{R}^{6\times18}$、伪逆运动力分配、$\ker(\mathbf{G})$ 的内力参数化。三臂共持一块板，要求板水平、内力在摩擦锥内；对比“平分运动力”与“$\mathbf{G}^{+}$ 分配”下各臂力矩峰值。
  \item \textbf{[跨章综合]} 内/外阻抗需要物体动力学参数（$\mathbf{M}_o$ 等）。论证：能否把分布式自适应搬运（无需负载惯性）的思想嫁接到内环，做“无需精确物体刚度的内力调节”？这与 \cref{sec:ma-arch} 的分布式架构有何关联？
\end{enumerate}
\end{exercise}

% ════════════════════════════════════════════════════════════════
\section{主从 vs 协调控制：一场关于体系结构的辩论}
\label{sec:ma-arch}

\begin{insight}[阶段小结：进入架构前的锚点]
到这里我们完成了协同操作的“数学 $+$ 控制”三连：(1) \cref{sec:ma-coopkin} 用对称表述把双臂运动拆成\emph{绝对}（物体运动）$+$\emph{相对}（内力通道）两个解耦坐标，核心对象是 $\mathbf{J}_{\text{abs}},\mathbf{J}_{\text{rel}}$；(2) \cref{sec:ma-duality} 用虚功原理证明速度约束 $\ker(\mathbf{J}_{\text{rel}})$ 与力约束 $\ker(\mathbf{G})$ 是同一闭链的对偶面孔，得出“内力 $=$ 相对运动的力共轭、必须闭环”；(3) \cref{sec:ma-internalforce} 把这个对偶变成 Caccavale 内/外阻抗双层控制器，并推广到 $N$ 臂 $\mathbf{h}=\mathbf{G}^{+}\mathbf{F}+\mathbf{W}\boldsymbol{\lambda}$。但前三节有一个未言明的假设：所有这些计算由“谁”来做、各臂“谁说了算”？我们默认了某个集中求解器同时算两臂——这正是接下来要审视的体系结构问题。
\end{insight}

\paragraph{动机：同一任务，两种“指挥结构”。}
玻璃要从 A 点移到 B 点。\textbf{方式一（主从）}：指定右手为“主臂”（master），给它轨迹 $\mathbf{T}_R^{\text{des}}(t)$ 去执行；左手是“从臂”（slave），任务是“无论主臂在哪，我都保持和主臂的固定相对位姿”。控制结构单向：主 $\to$ 从。\textbf{方式二（协调）}：不指定主从，给“物体”一条轨迹 $\mathbf{T}_{\text{abs}}^{\text{des}}(t)$（\cref{sec:ma-coopkin} 的绝对变量），两臂\emph{对称地}协商各自该怎么动，共同实现物体运动同时维持内力。控制结构对称。这两种方式不是实现细节的差别，而是\textbf{体系结构的根本分歧}，直接对应集中/分布/去中心化三分（\cref{tab:ma-arch-map}）。

\begin{table}[H]\centering\small
\caption{双臂控制方式与三大架构的对应}
\label{tab:ma-arch-map}
\begin{tabular}{llll}
\toprule
双臂控制方式 & 架构 & 谁拥有“全局意图” & 信息流 \\
\midrule
主从（master--slave） & 集中式（主臂是中心） & 主臂独占 & 单向 主 $\to$ 从 \\
对称协调（non-master--slave） & 分布式 & 物体目标共享给两臂 & 对称双向 \\
leader--follower 一致性 & 去中心化（leader 只是 1 个特殊节点） & leader 持有，经一致性传播 & 沿通信图传播 \\
\bottomrule
\end{tabular}
\end{table}

\subsection{主从控制：简单，但代价是什么}
\label{sec:ma-masterslave}

主从是历史最早、实现最简单的方案：\emph{只需为主臂规划轨迹，从臂目标由“主臂位姿 $+$ 固定相对变换”自动算出}，省掉了“为物体规划再分解到两臂”的步骤。早期遥操作双臂（一只手由人操控、另一只跟随）天然是主从的。但主从有一个根植于其结构的代价，正是 Uchiyama--Dauchez 极力反对它的原因\cite{uchiyama1987symmetric}\pz{存疑: 该对称方案权威年份为 ICRA 1988(源稿标 1987), 见 \cref{tab:ma-history}}。

\begin{insight}[主从的结构性缺陷]
主从控制把“内力管理”隐式地、不可控地推给了从臂的跟踪误差。主臂只管自己的位姿、完全不知道（也不关心）内力；从臂努力跟随主臂，但它的跟踪误差直接表现为相对位姿偏差，而相对位姿偏差 $=$ 内力（\cref{sec:ma-duality} 对偶）。于是\emph{内力变成了“从臂跟踪误差”的副产品，没有任何环节显式地测量和调控它}。主臂走得快一点、从臂跟不上，内力就飙升；从臂过冲，内力就反向。内力在主从结构里是“失控的自由变量”。
\end{insight}

用集中式的语言：主臂是单点，继承了集中式的所有优缺点——简单、有明确的全局协调者；但也有单点的脆弱（主臂的奇异、限位、故障整个系统跟着遭殃），且\emph{没有为“内力”这个协同操作核心量预留控制接口}。反事实推理：主从控制下若主臂经过奇异位形会怎样？主臂接近奇异 $\to$ 为实现期望末端速度需要巨大关节速度 $\to$ 实际运动滞后/抖动；从臂忠实跟随“主臂实际位姿”（含滞后/抖动），但从臂自己未必在奇异 $\to$ 它能精确复现主臂的抖动 $\to$ 两臂相对位姿剧烈波动 $\to$ 内力剧烈震荡 $\to$ 物体被反复挤压拉扯、可能脱手或损坏。对称协调下：物体轨迹是任务变量，两臂共同分担，单臂奇异可由冗余 $+$ 另一臂补偿，内力由内环独立守护、不被某臂奇异绑架。

\subsection{对称协调：把“物体”提升为一等公民}
\label{sec:ma-coord}

对称协调（Uchiyama--Dauchez non-master--slave）的核心，正是 \cref{sec:ma-coopkin} 的绝对/相对分解\cite{uchiyama1993symmetric}：共同目标是“物体轨迹”$\mathbf{T}_{\text{abs}}^{\text{des}}$（“分布式”的标志：目标是共享集体量、不属于任何单一 agent）；两臂\emph{对称地}实现它（绝对通道共同驱动物体，相对通道由内环独立守护内力）；内力是\emph{显式控制量}（有专门内环，不再是跟踪误差副产品）。对称协调相对主从的优势，恰是分布式相对集中式的优势（\cref{tab:ma-coord-vs}）。

\begin{table}[H]\centering\small
\caption{主从（集中式）vs 对称协调（分布式）}
\label{tab:ma-coord-vs}
\begin{tabular}{lll}
\toprule
维度 & 主从（集中式） & 对称协调（分布式） \\
\midrule
内力 & 失控副产品 & 显式控制量（内环） \\
单臂奇异/故障 & 拖累全局 & 冗余 $+$ 另一臂补偿 \\
任务语义 & “主臂去哪” & “物体去哪”（更贴合意图） \\
异构臂 & 主从角色僵化 & 可按能力分配绝对/相对权重（\cref{sec:ma-symmetry} 的 $\alpha$） \\
实现复杂度 & 低 & 中（需绝对/相对分解 $+$ 双环） \\
扩展到 $N>2$ & 困难（一主多从协调差） & 自然（绝对 $=$ 共同目标，内力 $=\ker(\mathbf{G})$） \\
\bottomrule
\end{tabular}
\end{table}

\begin{insight}[主从 $\to$ 对称协调与“集中 $\to$ 分布”是同一回事]
集中式有一个“什么都知道的中心”（主臂/中央 MPC），简单但脆弱、不可扩展；分布式把“全局意图”变成共享的集体目标（物体轨迹/团队目标），各 agent 对称协商，鲁棒且可扩展，代价是需要显式的协调机制（绝对/相对分解 / 共识 / ADMM）。\textbf{双臂的“主从 vs 协调”是多机“集中 vs 分布”在 $N=2$ 上的最小实例}——这正是把双臂放进多机协作课程的深层理由。
\end{insight}

\subsection{leader--follower 一致性：主从的“分布式软化”}
\label{sec:ma-leaderfollower}

主从（刚性，集中式）和对称协调（完全对等，分布式）之间，还有一个重要中间形态，它把共识（\cref{ch:mr-consensus}）直接用上：\textbf{leader--follower 一致性}。思想：保留一个 leader（持有任务意图，如物体目标轨迹），但 follower 不再“刚性跟随 leader 的瞬时位姿”，而是\emph{通过一致性协议逐步收敛}到与 leader 一致的运动。经典主从是 follower 直接复现 leader 位姿（含所有抖动/滞后），单向硬耦合；leader--follower 一致性是 follower 维护自己对“物体该怎么动”的估计 $\hat{\mathbf{x}}_i$，通过邻居通信（含从 leader 处）做一致性更新，驱动估计误差收敛到任意小：
\begin{equation}
  \dot{\hat{\mathbf{x}}}_i=-\sum_{j\in\mathcal{N}_i}a_{ij}(\hat{\mathbf{x}}_i-\hat{\mathbf{x}}_j)-b_i(\hat{\mathbf{x}}_i-\mathbf{x}_{\text{leader}}).
  \label{eq:ma-leaderfollow}
\end{equation}
这正是 leader-following consensus（带 leader 牵引项 $b_i$ 的 $\dot{\mathbf{x}}=-\mathbf{L}\mathbf{x}$ 变体，见 \cref{ch:mr-consensus}）。其收敛速率由通信图的代数连通性 $\lambda_2$ 决定，收敛所需“信息传播”正比于图直径——前置自测“最大值共识”的直觉在这里复用。为什么这是“软化”？见 \cref{tab:ma-leaderfollow-vs}。

\begin{table}[H]\centering\small
\caption{经典主从 vs leader--follower 一致性}
\label{tab:ma-leaderfollow-vs}
\begin{tabular}{lll}
\toprule
性质 & 经典主从 & leader--follower 一致性 \\
\midrule
follower 抖动放大 & 是（硬复现 leader） & 否（一致性滤波平滑） \\
扩展到 $N$ follower & 一主多从协调差 & 自然（多 follower 共识） \\
通信拓扑 & 必须 follower 直连 master & 任意连通图即可（信息可多跳传播） \\
鲁棒性 & leader 故障全崩 & 可设计 leader 切换 / 多 leader \\
架构归类 & 集中式 & 去中心化 \\
\bottomrule
\end{tabular}
\end{table}

\begin{insight}[双臂 leader--follower 一致性 vs 乐队的指挥与乐手（带边界）]
\textbf{像的部分}：指挥（leader）持有整体意图（曲目/物体轨迹），乐手（follower）不是机械复制指挥的每个手部抖动，而是把指挥的拍子作为参考、结合相邻乐手的演奏，平滑地协调到一致（一致性滤波）。\textbf{不像的部分}：乐手之间靠听觉（高带宽、低延迟）持续微调，机器人受 DDS 通信延迟与频率约束；乐队的“内力”（声部平衡）靠艺术经验，机器人靠 \cref{sec:ma-internalforce} 的显式内力环。\textbf{不要延伸到}：“既然乐队能无指挥即兴（完全去中心化），双臂也总能无 leader”——完全对称协调要求两臂共享物体目标且通信良好；在通信受限或任务天然有主次（一只手扶、一只手做精细操作）时，保留 leader 反而更合适，这就引出 \cref{sec:ma-symmetry} 的非对称。
\end{insight}

\subsection{三种架构的统一视角与选型}
\label{sec:ma-arch-select}

把三种架构放在统一框架里给出选型决策（\cref{tab:ma-arch-select}）。

\begin{table}[H]\centering\small
\caption{三种架构的选型决策}
\label{tab:ma-arch-select}
\begin{tabular}{lll}
\toprule
场景 & 推荐架构 & 理由 \\
\midrule
人遥操作一只手、另一只跟随 & 主从 & leader 是人，天然单向，内力靠人的力觉 \\
两工业臂固定协作、追求内力精度 & 对称协调 & 内力是质量关键，需显式控制 \\
$N$ 台移动机械臂、通信受限、可能掉线 & leader--follower 一致性 & 可扩展、容忍拓扑变化、抖动滤波 \\
一只手扶持、一只手精细操作 & 非对称协调（\cref{sec:ma-symmetry}） & 任务天然有主次，但要显式内力 \\
异构臂（能力差异大） & 对称协调 $+$ 加权分配 & 用 $\alpha$/加权伪逆把重活给强臂 \\
\bottomrule
\end{tabular}
\end{table}

\begin{insight}[通信受限场景偏好去中心化的一致性架构]
反事实推理：在“$N$ 台移动机械臂、通信会掉线”的场景硬上对称协调（要求全员共享物体目标、实时双向通信）会怎样？通信掉线时部分 agent 收不到物体目标更新，它们的绝对通道失去参考 $\to$ 各自漂移 $\to$ 内力失控。leader--follower 一致性更鲁棒：即使某 follower 暂时只能和邻居（非 leader）通信，一致性协议仍能通过多跳把 leader 意图传过去（代价是收敛慢一点，由 $\lambda_2$ 决定），不会立刻失控。这就是为什么通信受限的多机场景偏好去中心化的一致性架构——它把“对全局信息的依赖”降到“对邻居信息的依赖”。
\end{insight}

\begin{pitfall}[思维陷阱：认为“主从更简单所以更可靠”]
新手想法：“主从结构简单、代码少，工业上应该首选。”实际上：主从的简单是表面的——它把内力管理这个协同操作的核心难题“扫到地毯下”（变成从臂跟踪误差的副产品），并没有真正解决。在对内力敏感的任务（搬易碎物、精密装配）里，主从因内力失控反而最不可靠。简单 $\neq$ 可靠。正确理解：主从适合内力不敏感、有天然 leader（如人遥操作）的场景；对内力敏感的任务，多花的协调控制复杂度是值得的。
\end{pitfall}

\begin{pitfall}[概念误区：把 leader--follower 一致性等同于经典主从]
新手想法：“不都是有个 leader 让别人跟随吗？”实际上：经典主从是 follower 硬复现 leader 瞬时位姿（单向、放大抖动、要求直连）；leader--follower 一致性是 follower 通过共识协议逐步收敛到与 leader 一致（多跳传播、滤波平滑、容忍拓扑变化）。后者去中心化、前者集中式，鲁棒性和扩展性天差地别。正确理解：一致性把“刚性跟随”换成“协商收敛”，这是共识（\cref{ch:mr-consensus}）的直接应用。
\end{pitfall}

\begin{pitfall}[编程陷阱：对称协调里两臂的物体目标轨迹时钟不同步]
错误做法：两臂控制器各自从本地时钟生成 $\mathbf{x}_{\text{abs}}^{\text{des}}(t)$，时钟有偏移。现象：两臂对“物体此刻该在哪”理解不一致 $\to$ 相当于给相对通道注入误差 $\to$ 内力周期性波动。根本原因：对称协调要求两臂共享同一个物体目标，时钟不同步破坏了“共享”。正确做法：物体目标轨迹由统一时钟分发（或用一致性同步时钟），确保两臂在同一时刻引用同一物体目标。
\end{pitfall}

\begin{exercise}
\begin{enumerate}
  \item \textbf{[分析]} 复现“主臂经过奇异位形”：让主臂轨迹经过一个近奇异构型，记录主从控制下的内力峰值；再换对称协调，对比内力峰值（预测：主从显著更高）。
  \item \textbf{[设计]} 为“3 台移动机械臂协同抬桌、其中 1 台可能临时掉线”设计 leader--follower 一致性架构：写出每个 agent 的一致性更新律（含 leader 牵引项），说明掉线时系统如何降级而不崩溃，估计收敛速率与 $\lambda_2$ 的关系。
  \item \textbf{[跨章综合]} 论证：双臂的“主从 $\to$ leader--follower 一致性 $\to$ 对称协调”演进，与多机的“集中式 $\to$ 去中心化 $\to$ 分布式”是同一条主线在 $N=2$ 上的实例。指出三处精确对应和一处不对应（提示：双臂有强物理耦合即内力，一般多机协调如编队耦合更弱）。
  \item \textbf{[思考]} 主从控制把内力变成“失控副产品”，但遥操作双臂（人控主臂）却广泛用主从且工作良好。为什么？（提示：人通过主臂的力觉反馈隐式感知并调节内力——人脑充当了缺失的内力环。这说明“内力环”可由人提供，也说明全自主双臂为什么更需要显式内力控制。）
\end{enumerate}
\end{exercise}

% ════════════════════════════════════════════════════════════════
\section{对称与非对称协调}
\label{sec:ma-symmetry}

\paragraph{动机：现实中大多数双手任务是非对称的。}
\cref{sec:ma-arch} 的“主从 vs 对称协调”是\emph{架构}层面的对立（谁主导）。本节讨论一个正交维度：\emph{任务}本身是对称还是非对称？观察你自己的双手做事：拧瓶盖（一手固定瓶身、一手旋转盖）、穿针（一手持针、一手穿线）、切菜（一手按食材、一手运刀）、写字时另一只手按纸——\emph{绝大多数日常双手任务，两只手扮演不同角色}：一只通常是“稳定/参考”角色（变化慢、提供支撑），另一只是“操作/精细”角色（变化快、执行主要动作）。完全对称的任务（两人抬同一根杆、搓手）反而是少数。然而 \cref{sec:ma-coopkin} 的标准对称表述（绝对 $=$ 中点、相对 $=$ 右减左）默认两臂对等——直接用它处理非对称任务会“水土不服”：以“左手扶碗、右手搅拌”为例，对称表述把“绝对运动”定义为两手中点的运动，但搅拌任务里我们关心的是“碗（左手）放在哪里、保持稳定”（左手的绝对运动）与“勺子（右手）相对碗怎么转圈”（相对运动），而“两手中点”这个量在搅拌里既不是碗的位置也不是勺的位置，是个尴尬的几何中点；给“中点”规划轨迹需要把“碗稳定 $+$ 勺转圈”反向折算成中点轨迹，绕一大圈且任务描述被扭曲。

\begin{insight}[对称表述的“绝对 $=$ 中点”是关于谁代表集体运动的民主假设]
它对真正对称的任务（两手地位等价）是最优的，但对非对称任务，“集体运动”更应该由\emph{承担参考角色的那只手}代表。任务的对称性应当决定“绝对变量怎么定义”——这就是 \cref{sec:ma-coopkin} 引入连续参数 $\alpha$ 的全部意义。
\end{insight}

\subsection{用 $\alpha$ 统一对称与非对称}
\label{sec:ma-alpha-unify}

回顾 \cref{eq:ma-alpha} 的加权绝对 twist $\boldsymbol{\mathcal{V}}_{\text{abs}}^{(\alpha)}=\alpha\,\bar{\boldsymbol{\mathcal{V}}}_L+(1-\alpha)\,\bar{\boldsymbol{\mathcal{V}}}_R$。$\alpha$ 是一个把“集体运动以谁为代表”参数化的连续旋钮（\cref{tab:ma-alpha-table}）。

\begin{table}[H]\centering\small
\caption{$\alpha$ 把对称/非对称统一进一个连续参数}
\label{tab:ma-alpha-table}
\begin{tabular}{llll}
\toprule
$\alpha$ & 绝对变量代表 & 对应任务类型 & 例子 \\
\midrule
$0.5$ & 两手中点（民主） & 对称任务 & 两人抬杆、搓手 \\
$1.0$ & 左手（左手为参考） & 非对称，左扶右做 & 左手扶碗、右手搅拌 \\
$0.0$ & 右手（右手为参考） & 非对称，右扶左做 & 右手持工件、左手加工 \\
$\in(0,1)$ & 加权（按能力/惯量） & 异构臂 & 强臂权重高 \\
\bottomrule
\end{tabular}
\end{table}

\noindent\textbf{$\alpha$ 与架构的关系（关键）}：$\alpha$ 是\emph{任务}维度（任务对称性），\cref{sec:ma-arch} 的主从/对称是\emph{架构}维度（控制组织）。两者正交但相关：对称任务（$\alpha=0.5$）$+$ 对称架构 $=$ 经典 Uchiyama--Dauchez；非对称任务（$\alpha=1$）天然倾向于让“参考手（扶持手）”承担类 leader 角色——但\emph{这不等于主从}！即使 $\alpha=1$（以左手为绝对参考），内力仍由独立内环显式控制（\cref{sec:ma-internalforce}），左手并不“无视内力”。这是非对称协调与主从的关键区别：非对称协调保留了内力的显式控制，主从没有。

\begin{insight}[$\alpha$ 把“任务非对称性”和“架构主从性”解耦了]
一个常见混淆是把“一只手扶、一只手做”直接等同于“主从控制”——错。前者是任务非对称（用 $\alpha$ 表达），后者是架构集中式（用单向信息流表达）。你完全可以用\emph{对称架构（两臂协商、内力显式）执行非对称任务（$\alpha=1$）}：扶持手当绝对参考，操作手做相对运动，但两臂仍对等地协商、内力仍被独立守护。$\alpha=1$ 选的是“谁代表物体运动”，不是“谁单向指挥谁”。
\end{insight}

\subsection{扩展相对雅可比的落地}
\label{sec:ma-extjac-landing}

\cref{sec:ma-coopkin} 给出了非对称任务的扩展相对雅可比 $\mathbf{J}_{\text{ext}}=[\mathbf{J}_{\text{rel}};\mathbf{J}_{\text{abs}}^{(k)}]$（\cref{eq:ma-extjac}）。对“左手扶碗、右手搅拌”（$\alpha=1$，左手为绝对参考）：$\boldsymbol{\mathcal{V}}_{\text{abs}}^{(1)}=\bar{\boldsymbol{\mathcal{V}}}_L$ 直接控制左手（碗）的运动——任务描述清晰（碗放哪、稳不稳）；$\boldsymbol{\mathcal{V}}_{\text{rel}}=\bar{\boldsymbol{\mathcal{V}}}_R-\bar{\boldsymbol{\mathcal{V}}}_L$ 控制勺相对碗的运动——直接就是“搅拌动作”。控制律拆成两个解耦任务（任务优先级写法）：优先级 1（高）左手绝对运动 $=$ 碗的期望运动（保持碗稳定/到位），在 $\ker(\mathbf{J}_{\text{rel}})$ 中与碗稳定相容的子空间实现；优先级 2 右手相对左手运动 $=$ 搅拌轨迹；内力环（并行）相对位姿偏差的法向分量 $\to$ 维持勺压碗底的接触力；冗余 $\to$ 自碰撞回避、关节限位。

\begin{note}
\textbf{非对称内力常是单向接触力.}\;对称任务的内力是“两手对夹”（双向挤压）；非对称任务的内力常是“操作手对参考物的单向作用力”（如勺压碗底、刀压砧板、笔压纸）。这在 grasp matrix 框架里对应不同的内力子空间结构——对称是 $\ker(\mathbf{G})$ 的“反对称”部分，非对称更关注“操作手 $\to$ 参考”方向的法向力。处理非对称任务时，把内力期望值设计成“维持必要接触力”而非“对夹力”，即在 $\ker(\mathbf{G})$ 里选对应子空间。
\end{note}

\subsection{把非对称协调连到任务分配}
\label{sec:ma-roleassign}

\cref{sec:ma-perspective} 练习 3 已埋下伏笔：双臂的“角色分配（哪只手当参考、哪只手当操作）”是任务分配的退化形式。任务分配把“$N$ 机器人 $\times$ $M$ 任务，谁干哪件最划算”形式化为指派问题；双臂非对称任务里也有一个角色分配：\emph{2 只手 $\times$ 2 个角色（参考/操作），谁当参考谁当操作}？（\cref{tab:ma-roleassign}）。

\begin{table}[H]\centering\small
\caption{一般任务分配 vs 双臂角色分配}
\label{tab:ma-roleassign}
\begin{tabular}{lll}
\toprule
维度 & 一般任务分配 & 双臂角色分配（本节） \\
\midrule
agent 数 & $N$（大） & 2 \\
任务/角色数 & $M$ & 2（参考、操作） \\
代价 & 距离/时间 & 哪只手当参考更稳（工作空间、惯量、当前构型） \\
关键约束 & agent 可独立执行任务 & \textbf{两角色强物理耦合}，不能独立 \\
求解 & 匈牙利/拍卖/CBBA & 简单比较（$2!=2$ 种）$+$ 物理判据 \\
\bottomrule
\end{tabular}
\end{table}

\noindent\textbf{关键区别}：一般任务分配假设 agent 能独立执行各自任务（解耦或弱耦合）；双臂角色分配里两个角色\emph{通过共持物体强耦合}——参考手和操作手的运动学、内力完全交织。所以双臂角色分配不能照搬“独立指派”求解，而要结合 \cref{sec:ma-coopkin,sec:ma-internalforce} 的耦合约束。但“为操作选合适的角色分工”这个\emph{思想}是相通的：都是把“谁做什么”显式化为决策变量。反事实推理：如果角色分配错了（让工作空间更受限的手当操作手）会怎样？操作手要执行精细动作（搅拌/拧/切），需要较大的灵活工作空间和远离奇异的构型；若把它分给一只快到限位、临近奇异的手，操作精度和范围都受损，甚至中途卡死。正确的角色分配应把“操作”角色给当前构型更灵活、操作度更高的手——这是一个（退化的）能力匹配问题。在长序列任务里，角色甚至需要中途切换（regrasp/换手），这又回到 TAMP（离散角色决策 $+$ 连续运动）。

\begin{insight}[对称/非对称选型决策树]
两手是否扮演相同角色（可互换）？\emph{是} $\to$ 对称任务（$\alpha=0.5$，绝对 $=$ 中点，内力 $=$ 对夹力，例：两人抬杆、对称装配）；\emph{否} $\to$ 非对称任务，哪只手是“参考/稳定”角色就设它为绝对参考（$\alpha=1$ 或 $0$），用扩展相对雅可比控制 $[\text{相对运动}+\text{参考手绝对运动}]$，内力 $=$ 操作手对参考的接触力（常单向，例：扶切、持拧、按写）。\emph{正交地}，无论对称/非对称都要选架构（\cref{sec:ma-arch}）：主从 / leader--follower 一致性 / 对称协调——对内力敏感就别用纯主从。
\end{insight}

\begin{pitfall}[概念误区：把“非对称任务”等同于“主从控制”]
新手想法：“一只手扶、一只手做，那不就是扶的当主、做的当从（主从控制）吗？”实际上：非对称是任务属性（两手角色不同，用 $\alpha$ 表达“谁代表物体运动”），主从是架构属性（单向信息流、内力失控）。你可以用对称架构（两臂协商、内力显式环）执行非对称任务——扶持手当绝对参考，但内力仍被独立守护。混淆二者会让你在需要内力精度的非对称任务（如精密装配中一手扶一手插）里错误地用主从，导致内力失控。正确理解：任务对称性（$\alpha$）和架构主从性是两个正交维度，分别决策。
\end{pitfall}

\begin{pitfall}[思维陷阱：对称表述（绝对 $=$ 中点）适用于所有任务]
新手想法：“Uchiyama--Dauchez 对称表述是标准，所有双臂任务都用它。”实际上：对称表述对真正对称的任务最优，但对非对称任务，“两手中点”缺乏清晰物理意义，强行使用会扭曲任务描述。正确思维：用 $\alpha$ 匹配任务对称性——对称任务 $\alpha=0.5$，非对称任务 $\alpha=1/0$（以参考手为绝对变量）。
\end{pitfall}

\begin{pitfall}[编程陷阱：非对称任务里沿用“对夹内力”模型]
错误做法：扶切任务（刀压砧板）里，把内力建模成两手对夹（双向挤压）。现象：控制器试图让“砧板和刀互相对夹”，但任务只需要“刀向下压砧板”的单向接触力，对夹模型产生多余的水平内力，刀打滑。根本原因：非对称任务的内力常是单向接触力（操作手 $\to$ 参考），不是对夹。正确做法：按任务定义内力方向——对称任务用对夹，非对称用单向接触力，在 $\ker(\mathbf{G})$ 里选对应子空间。
\end{pitfall}

\begin{exercise}
\begin{enumerate}
  \item \textbf{[标注]} 对以下任务标注对称性（给出建议 $\alpha$）和内力类型（对夹/单向接触）：(1) 两人抬沙发；(2) 拧瓶盖；(3) 穿针引线；(4) 双手系鞋带；(5) 一手扶梯一手钉钉子。说明每个的“绝对参考”该选谁。
  \item \textbf{[实现]} 用扩展相对雅可比实现“左手持板、右手在板上画直线”：$\alpha=1$，优先级 1 保持板稳定，优先级 2 让右手沿板面直线运动，内力环维持右手笔尖对板的法向压力 $5$\,N。
  \item \textbf{[设计]} 为长序列任务“把零件从左手递到右手再装配”设计角色切换：递交前左手为操作手、右手为参考；递交后角色互换。说明切换时刻如何避免内力突变（提示：切换瞬间 $\alpha$ 不能跳变，需平滑过渡）。
  \item \textbf{[跨章综合]} 论证：“长序列双臂操作中的角色分配 $+$ 换手时机”是一个 TAMP 问题（离散角色决策 $+$ 连续运动），而单步的角色选择是任务分配的退化（$N=2$、强耦合）。给出一个需要换手的任务并说明 TAMP 如何决策换手点。
\end{enumerate}
\end{exercise}

% ════════════════════════════════════════════════════════════════
\section{多臂负载分配：$N\ge 3$ 的优化与分布式}
\label{sec:ma-allocation}

\paragraph{动机：为什么 $N\ge 3$ 后“平分”不再可行。}
双臂 $N=2$ 时，运动力平分（每手 $\tfrac12\mathbf{F}$）$+$ 内力对夹，简单且常常够用——因为 $N=2$ 的内力空间只有 $6$ 维，自由度小，对称情形下平分恰好最优。$N\ge3$ 后情况质变（回顾 \cref{sec:ma-perspective} 反面教材“四臂抬钢板”）：三臂抬一块质心偏置的板，臂 1、臂 2 离质心远（力臂大）、臂 3 离质心近（力臂小），三臂电机扭矩上限还各不相同。朴素“三等分”每臂承担 $1/3$ 重量——但臂 3 离质心近、力臂小，承担 $1/3$ 重量需要的关节力矩可能反而最大（要在小力臂上产生大力），或者臂 3 电机最弱；谁先到力矩上限谁就饱和，份额甩给邻居引发连锁过载。

\begin{insight}[$N\ge3$ 时负载分配从“平分公式”变成“带约束优化”]
存在 $6N-6$ 维内力自由度可以利用。这个自由度既是麻烦（不再有唯一答案）也是机会（可以优化——让强臂多担、弱臂少担，留足防滑裕度，远离力矩饱和）。\textbf{双臂的“平分”只是这个优化在 $N=2$、对称、无能力差异时的退化最优解}。一旦异构或 $N\ge3$，必须真正去优化。
\end{insight}

\subsection{集中式 QP 形式}
\label{sec:ma-qp}

把负载分配写成标准 QP。决策变量是各手 wrench 堆叠 $\mathbf{h}=[\mathbf{h}_1;\dots;\mathbf{h}_N]\in\mathbb{R}^{6N}$。\textbf{等式约束（必须实现物体运动）}：各手合成 wrench 等于物体所需 wrench（含重力、惯性）$\mathbf{G}\mathbf{h}=\mathbf{F}_{\text{load}}^{\text{cmd}}$——这把 $\mathbf{h}$ 限制在仿射子空间 $\mathbf{G}^{+}\mathbf{F}_{\text{load}}^{\text{cmd}}+\ker(\mathbf{G})$，剩下 $6N-6$ 维内力自由度待优化。\textbf{不等式约束（物理可行）}：摩擦锥（接触不打滑）——每个接触力 $\mathbf{f}_i$ 落在摩擦锥内 $\|\mathbf{f}_i^{\text{tangent}}\|\le\mu_i f_i^{\text{normal}}$，且 $f_i^{\text{normal}}\ge f_{\min}$（最小法向力防滑），摩擦锥是二阶锥，可用金字塔近似线性化为线性不等式；能力上限（电机不饱和）——各臂关节力矩 $\boldsymbol{\tau}_i=\mathbf{J}_i^\top\mathbf{h}_i$ 满足 $|\boldsymbol{\tau}_i|\le\boldsymbol{\tau}_{\max,i}$。\textbf{代价函数}（几选其一或加权，\cref{tab:ma-qp-cost}）。综合 QP：
\begin{equation}
  \boxed{\;\min_{\mathbf{h}}\ \tfrac12\mathbf{h}^\top\mathbf{Q}\mathbf{h}+\mathbf{c}^\top\mathbf{h}
  \quad\text{s.t.}\quad \mathbf{G}\mathbf{h}=\mathbf{F}_{\text{load}}^{\text{cmd}},\ \ \mathbf{C}\mathbf{h}\le\mathbf{d}\;}
  \label{eq:ma-qp}
\end{equation}
$\mathbf{Q}$ 编码代价（加权范数 $+$ 内力正则），$\mathbf{C}\mathbf{h}\le\mathbf{d}$ 堆叠摩擦锥线性化 $+$ 力矩上限。这是一个凸 QP，可用现成 QP 求解器实时求解。

\begin{table}[H]\centering\small
\caption{多臂负载分配的几种代价目标}
\label{tab:ma-qp-cost}
\begin{tabular}{lll}
\toprule
目标 & 代价 & 含义 \\
\midrule
最小接触力范数 & $\sum_i\|\mathbf{h}_i\|^2$ & 最省力，但可能不公平 \\
最小化最大归一化力 & $\min\max_i\|\mathbf{h}_i\|/\tau_{\max,i}$（min--max，化为 LP/QP） & 公平，远离饱和，异构友好 \\
加权范数 & $\sum_i w_i\|\mathbf{h}_i\|^2$，$w_i\propto1/\text{能力}_i$ & 强臂多担 \\
内力正则 & $\|\mathbf{f}_{\text{int}}-\mathbf{f}_{\text{int}}^{\text{des}}\|^2$ & 把内力拉向防滑裕度区间 \\
\bottomrule
\end{tabular}
\end{table}

\begin{insight}[QP 与协同搬运的“参数化”是同一回事的两种表述]
多臂负载分配的 QP 和协同搬运章的 $\mathbf{f}_{\text{c}}=\mathbf{G}^{+}\mathbf{F}_{\text{load}}+(\mathbf{I}-\mathbf{G}^{+}\mathbf{G})\mathbf{f}_{\text{int}}$ 是同一回事的两种表述——参数化（$\mathbf{f}_{\text{int}}$ 自由）vs 优化（在约束下选最优 $\mathbf{f}_{\text{int}}$）。协同搬运章给的是“内力是自由参数”的结构性认识，本节把它升级为“在摩擦锥 $+$ 能力约束下选最优内力”的可执行优化。从结构到优化，正是从“理解”到“工程”的一步。
\end{insight}

\textbf{平分 vs QP 的数值对比（定性预期）}。对“三臂抬偏置板、臂 3 能力弱（$\tau_{\max,3}$ 只有一半）”的场景，关键差异在“最弱臂的归一化负载”$\|\mathbf{h}_i\|/\tau_{\max,i}$——谁先到 $1$ 谁就饱和（\cref{tab:ma-equal-vs-qp}）。

\begin{table}[H]\centering\small
\caption{朴素平分 vs 加权 QP（三臂、臂 3 能力弱）}
\label{tab:ma-equal-vs-qp}
\begin{tabular}{lll}
\toprule
指标 & 朴素平分（各 $1/3$） & 加权 QP（$w_3$ 高 $\to$ 臂 3 少担） \\
\midrule
臂 3 承担力占比 & $\sim33\%$（不顾它能力弱） & 显著 $<33\%$（重活转给臂 1、臂 2） \\
臂 3 归一化负载 $\|\mathbf{h}_3\|/\tau_{\max,3}$ & 最高，\emph{最先逼近饱和} & 被压低，远离饱和 \\
系统“最大归一化负载” & 由最弱臂决定，偏高 & 被均衡，整体下降 \\
防滑/摩擦锥 & 不保证（无 $f_{\min}$ 约束时） & 显式约束，保证 \\
\bottomrule
\end{tabular}
\end{table}

\begin{insight}[优化目标应对准瓶颈：最小化“最大归一化负载”]
平分优化的是“绝对力均等”，QP（加权 / min--max）优化的是“归一化负载均衡”——后者才是异构系统真正该关心的。直观判据：\emph{系统的搬运能力上限由“最先饱和的那只手”决定}，所以应该最小化“最大归一化负载”（min--max），而非让绝对力相等。这与任务分配里“最小化 makespan（最晚完成的 agent）”而非“总代价”的思想完全一致——多机系统的瓶颈往往是最差的那个个体，优化目标应对准瓶颈。
\end{insight}

\begin{insight}[$f_{\min}$ 约束是把“防滑”写进优化的方式]
反事实推理：QP 里如果去掉“最小法向力 $f_{\min}$”约束会怎样？求解器为了最小化代价，会把某些接触法向力压到接近零——数学上最省力，但物理上意味着那只手“几乎没夹住”，稍有扰动就打滑脱落。$f_{\min}$ 约束（以及内力正则把内力拉向防滑裕度）是把“防滑”这个物理必需写进优化的方式。没有它，最优解在数学上完美、在物理上危险——这是优化驱动控制的典型陷阱：目标函数没编码的约束，求解器不会替你考虑。
\end{insight}

\noindent\textbf{grasp matrix 与摩擦锥的实现要点}（小段伪代码，载体无关）。$\mathbf{G}$ 的每个 $6\times6$ 块按 \cref{eq:ma-grasp} 填（力 $\to$ 力为 $\mathbf{I}_3$、力 $\to$ 力矩为 $[\mathbf{r}_i]_\times$、力矩 $\to$ 力矩为 $\mathbf{I}_3$）；摩擦锥金字塔近似把 $\|\mathbf{f}^{\text{tangent}}\|\le\mu f^{\text{normal}}$ 写成 $|f_x|\le\mu f_z$、$|f_y|\le\mu f_z$（$4$ 条线性不等式，设接触法向为 $+z$）：用 \verb|build_grasp_matrix(contact_pts, com)| 组装 $\mathbf{G}$、用 \verb|solve_load_distribution(...)| 解 \cref{eq:ma-qp}，关键参数为按 $1/\tau_{\max,i}$ 设的权重 \verb|w| 与最小法向力 \verb|f_min|；解后用 \verb|G @ h - F_cmd| $\approx\mathbf{0}$ 验证合成正确。具体求解器调用属工程细节，从略。

\subsection{分布式版本：用 ADMM 把 QP 拆给各 agent}
\label{sec:ma-admm}

集中式 QP 需要一个中央节点收集所有 $\mathbf{r}_i,\tau_{\max,i}$、求解、下发 $\mathbf{h}_i$。但真实多机系统（$N$ 台移动机械臂、分布在多个计算节点）希望\emph{每个 agent 只用本地 $+$ 邻居信息}求解——这正是 ADMM（\cref{sec:cs-admm}）的用武之地。这里的耦合约束是 $\mathbf{G}\mathbf{h}=\mathbf{F}_{\text{load}}^{\text{cmd}}$（所有手的 wrench 必须合成出物体所需 wrench）——一个把所有 agent 绑在一起的全局等式约束。ADMM 分解思路：全局问题 $\min\sum_i J_i(\mathbf{h}_i)$ s.t. $\sum_i\mathbf{G}_i\mathbf{h}_i=\mathbf{F}_{\text{cmd}}$、$\mathbf{h}_i\in$ 局部可行集（摩擦锥/力矩）。引入“物体 wrench 残差”作为对偶变量（拉格朗日乘子 $\boldsymbol{\lambda}$，物理上是“物体当前合成 wrench 与目标的差”对应的协调价格），每个 agent $i$ 的本地子问题（只需自己的 $\mathbf{G}_i,\mathbf{J}_i,\tau_{\max,i}$ $+$ 全局 $\boldsymbol{\lambda}$）为
\begin{equation}
  \mathbf{h}_i\leftarrow\arg\min_{\mathbf{h}_i}\ J_i(\mathbf{h}_i)+\boldsymbol{\lambda}^\top(\mathbf{G}_i\mathbf{h}_i)+\tfrac{\rho}{2}\|\mathbf{G}_i\mathbf{h}_i-\mathbf{z}_i\|^2
  \quad\text{s.t. }\mathbf{h}_i\in\text{局部可行集},
  \label{eq:ma-admm-local}
\end{equation}
$\mathbf{z}$ 更新（一致性，把 $\sum\mathbf{G}_i\mathbf{h}_i=\mathbf{F}_{\text{cmd}}$ 的残差均摊/协商），$\boldsymbol{\lambda}$ 更新（对偶上升，广播残差）$\boldsymbol{\lambda}\leftarrow\boldsymbol{\lambda}+\rho(\sum_i\mathbf{G}_i\mathbf{h}_i-\mathbf{F}_{\text{cmd}})$。每个 agent 只解一个\emph{小 QP}（自己的 $6$ 维 wrench $+$ 局部约束），全局协调通过交换 $\boldsymbol{\lambda}$（物体 wrench 残差，$6$ 维，很小）完成；$\boldsymbol{\lambda}$ 的交换可用共识（\cref{ch:mr-consensus}）：若没有中央节点广播，各 agent 通过邻居一致性估计全局残差。

\begin{insight}[$\boldsymbol{\lambda}$ 的物理意义是“物体 wrench 的协调价格”]
它告诉每个 agent“物体现在还差多少 wrench，你该多出还是少出力”。这与拍卖机制的“价格”、ADMM 的“对偶变量”是同一个思想：\emph{用一个共享的标量/向量价格协调分布式资源分配}，每个 agent 只需看价格（而非别人的全部状态）就能做出协调的本地决策。协同操作的力分配，至此完全融入分布式优化框架。
\end{insight}

\begin{insight}[ADMM 力分配的 $\boldsymbol{\lambda}$ vs 市场价格信号（带边界）]
\textbf{像的部分}：市场里没有中央计划者，每个买卖方只看价格就能做出协调决策，价格通过供需失衡（残差）自动调整；ADMM 里每个 agent 只看 $\boldsymbol{\lambda}$（物体 wrench 缺口的价格）就能决定本地出力，$\boldsymbol{\lambda}$ 通过残差 $\sum\mathbf{G}_i\mathbf{h}_i-\mathbf{F}_{\text{cmd}}$ 调整。\textbf{不像的部分}：市场价格是分散主体自利博弈的涌现，ADMM 的 $\boldsymbol{\lambda}$ 是为最小化全局代价而设计的更新律、有收敛保证；市场可能有信息不对称、欺诈，ADMM 假设 agent 诚实协作。\textbf{不要延伸到}：“既然市场无需协调就能均衡，多臂也无需任何通信”——ADMM 仍需交换 $\boldsymbol{\lambda}$（哪怕只是 $6$ 维、可经共识传播），完全零通信的力分配要靠隐式力反馈，是另一条更激进的路线。
\end{insight}

\begin{table}[H]\centering\small
\caption{力分配：集中式 QP vs 分布式 ADMM}
\label{tab:ma-qp-vs-admm}
\begin{tabular}{lll}
\toprule
考量 & 集中式 QP & ADMM 分布式 \\
\midrule
求解延迟 & 单次 QP（$6N$ 维），毫秒级，$N$ 大时变慢 & 各 agent 并行解小 QP，但需多轮 ADMM 迭代通信 \\
通信需求 & 中央收集全局、下发 $\mathbf{h}_i$ & 仅交换 $\boldsymbol{\lambda}$（$6$ 维），可经共识 \\
单点故障 & 中央节点挂 $=$ 全崩 & 无中央节点，鲁棒 \\
某 agent 饱和/掉线 & 重解 QP（去掉该 agent 的列） & 该 agent 退出，其余继续协商（份额自动转移） \\
适用规模 & $N$ 小到中（$2$--$6$） & $N$ 中到大、分布式部署 \\
\bottomrule
\end{tabular}
\end{table}

\begin{insight}[“集中 QP vs 分布式 ADMM”第三次重复本课程核心主题]
集中式简单但有单点、不可扩展；分布式（ADMM/共识）需要协调机制但鲁棒、可扩展。\cref{sec:ma-arch} 在控制架构上、本节在力分配上，都是这同一个“集中 vs 分布”权衡的具体化。当你理解了这一点，本章的“主从 vs 协调”“集中 QP vs ADMM”就不是孤立的知识点，而是同一个多机协作根本矛盾的不同切面。
\end{insight}

\begin{pitfall}[编程陷阱：QP 力分配里摩擦锥用二阶锥但喂给只支持线性约束的求解器]
错误做法：把 $\|\mathbf{f}^{\text{tangent}}\|\le\mu f^{\text{normal}}$（二阶锥）直接当线性约束塞进只解 QP $+$ 线性约束的求解器。现象：求解器报错或静默把锥当成别的约束，解出的力实际违反摩擦锥 $\to$ 打滑。根本原因：摩擦锥是二阶锥约束，纯 QP 求解器不直接支持；需要金字塔线性化（$4$--$8$ 条线性不等式近似锥）或用 SOCP 求解器。正确做法：要么金字塔线性化后用 QP（保守，锥被内接近似），要么用 SOCP 精确处理；线性化边数越多越接近真锥但 QP 越大。
\end{pitfall}

\begin{pitfall}[概念误区：认为“分布式 ADMM 一定比集中式 QP 慢”]
新手想法：“ADMM 要迭代很多轮还要通信，肯定比直接解一个 QP 慢。”实际上：取决于规模和部署。$N$ 小、单机时集中 QP 更快（一次求解）；但 $N$ 大、分布式部署时，集中式要把全局信息汇总到一点（通信瓶颈）$+$ 解一个大 QP（$O(N^3)$ 级），而 ADMM 各 agent 并行解小 QP，总 wall-clock 可能更短，且无单点瓶颈。正确理解：集中 vs 分布的选择取决于 $N$、是否分布式硬件、通信代价，不是单纯的“迭代 $=$ 慢”。
\end{pitfall}

\begin{pitfall}[思维陷阱：把 $N$ 臂内力当成“每对相邻臂各自调内力”]
新手想法：“$N$ 臂力分配 $=$ 让每对相邻的手各自维持对夹内力。”实际上：$N$ 臂内力是整体的 $6N-6$ 维子空间上的优化，各对内力相互耦合（一只手同时参与多对），且要全局满足“合成出物体 wrench”。局部地各调各的对夹，无法保证全局物体运动正确、也无法全局最优分配。正确思维：用全局 QP（或 ADMM 分解的全局一致性）在 $\ker(\mathbf{G})$ 整体上优化，而非局部成对处理。
\end{pitfall}

\begin{exercise}
\begin{enumerate}
  \item \textbf{[实现]} 跑通三臂 QP 力分配。改变臂 3 的位置（从靠近质心移到远离质心）和 $\tau_{\max,3}$，观察最优分配如何变化；验证：当某臂能力很弱时，QP 自动让它少担、邻居多担。
  \item \textbf{[推导]} 把摩擦锥的金字塔线性化写清楚：给定接触法向 $\mathbf{n}_i$ 和摩擦系数 $\mu_i$，用 $m$ 边金字塔近似摩擦锥，写出对应 $m$ 条线性不等式；讨论 $m$ 增大时的精度-计算量权衡。
  \item \textbf{[设计扩展·研究级]} 实现 ADMM 分布式力分配：三臂各跑本地 QP，通过交换对偶变量 $\boldsymbol{\lambda}$（物体 wrench 残差）协调；对比集中式 QP 的解（应收敛到一致），记录 ADMM 收敛所需轮数；再模拟“臂 2 掉线”，验证其余两臂能否重新协商出可行分配。
  \item \textbf{[跨章综合·研究级]} 把 ADMM、共识、本节力分配串成一个完整系统：$N$ 台移动机械臂协同搬运，物体目标用 leader--follower 一致性达成，力分配用本节 ADMM，且 $\boldsymbol{\lambda}$ 的全局残差通过共识估计（无中央节点）。画出系统框图，标出每个模块来自哪一章，并指出该系统在耦合强度谱（\cref{sec:ma-perspective}）上的位置。
\end{enumerate}
\end{exercise}

% ════════════════════════════════════════════════════════════════
\section{与单系统双臂内容的关系，及学习范式简述}
\label{sec:ma-boundary}

\paragraph{两套讲双臂的内容：互补而非重复。}
本项目有\emph{两套}讲双臂的内容，立场不同、互补：机械臂方向的双臂子课程站在\emph{单系统操作}立场，把双臂当一台 $14$ 自由度机器，深入到运动学/规划/力控的工程实现（库 API、奇异处理、约束规划、ROS2 部署）；本章站在\emph{多机协作}立场，把双臂当 $N=2$ 多机协同，给绝对/相对对偶、架构选型与 $N$ 臂推广的结构性认识。\cref{tab:ma-two-tracks} 是总分工表。

\begin{table}[H]\centering\small
\caption{单系统双臂内容 vs 本章多机视角的总分工}
\label{tab:ma-two-tracks}
\begin{tabular}{p{1.8cm}p{5.6cm}p{5.0cm}}
\toprule
维度 & 单系统视角（机械臂方向） & 本章（多机协作视角） \\
\midrule
立场 & 双臂 $=$ 一台 $14$ 自由度机器 & 双臂 $=N=2$ 多机协同的特例 \\
运动学 & $\mathbf{J}_{\text{aug}},\mathbf{J}_{\text{rel}}$ 完整推导 $+$ 工程实现、奇异/阻尼细节 & 从对称表述推导 $\mathbf{J}_{\text{rel}}$ 来历，引入 $\mathbf{J}_{\text{abs}}$ 与 $\alpha$，推广到 $N$ 臂（\cref{sec:ma-coopkin}） \\
规划 & 约束流形采样、TSR、CBiRRT、TAMP & 不涉及；只在角色分配处连到任务分配（\cref{sec:ma-symmetry}） \\
力控 & object/internal 阻抗工程、无源性、波变量、ROS2 力控 & 内力的 $\ker(\mathbf{G})$ 推导与 $N$ 臂推广、内/外解耦的体系结构含义（\cref{sec:ma-internalforce}） \\
架构 & 较少展开主从 vs 协调的多机辨析 & 主从/leader--follower/对称协调与三架构对应（\cref{sec:ma-arch}） \\
多臂 $N\ge3$ & 聚焦双臂 & grasp QP $+$ ADMM 分布式力分配（\cref{sec:ma-allocation}） \\
学习 & ALOHA/ACT/Diffusion Policy 深入 & 仅简述（下面） \\
力空间统一 & grasp matrix 简述 & $\ker(\mathbf{J}_{\text{rel}})/\ker(\mathbf{G})$ 对偶（\cref{sec:ma-duality}） \\
\bottomrule
\end{tabular}
\end{table}

\begin{insight}[同一对象值得从两个视角各讲一遍]
单系统视角给你“把这台双臂机器调好”的工程深度（API、奇异、ROS2），多机视角给你“把双臂知识长进多机系统”的结构性认识（对偶、架构、$N$ 臂推广）。一个工程师两者都需要：前者让你今天就能让双臂跑起来，后者让你明天能把它扩展成多机协作系统而不推倒重来。视角不是非此即彼，是同一座山的两条上山路。
\end{insight}

\paragraph{学习范式简述：从“建模控制”到“模仿学习”。}
$2022$ 年 ALOHA 之后\cite{zhao2023aloha}，双臂操作出现了一条与经典控制并行的\emph{学习范式}路线。本章主体讲经典控制（它可分析、可推广到多机），这里简述学习范式并说清两者关系。核心转变：经典控制要你\emph{显式写出}相对雅可比、内力分解、阻抗律；学习范式让机器\emph{从人类演示中学}双臂协调，不显式建模运动学/力学（\cref{tab:ma-learning}）。\rebuilt{下表年份、归属转述自源稿，待核。}

\begin{table}[H]\centering\small
\caption{双臂学习范式代表工作（多机视角注解）}
\label{tab:ma-learning}
\begin{tabular}{p{3.2cm}p{1.1cm}p{8.0cm}}
\toprule
工作 & 年份 & 核心思想（多机视角注解） \\
\midrule
ALOHA / ACT\cite{zhao2023aloha} & 2022 & 低成本遥操作采数据 $+$ Action Chunking Transformer。遥操作天然是\emph{主从}（人控 leader 臂、follower 复现），数据里隐含人脑提供的内力协调 \\
Mobile ALOHA & 2024 & 加移动底盘，$14$ 关节 $+$ 2D 底盘 $=16$D 动作；co-training 大幅提升成功率 \\
Diffusion Policy\cite{chi2023diffusion} & 2023 & DDPM 去噪生成动作 chunk，天然多模态、精度高\rebuilt{待修(勿原位重建): 原文(RSS 2023)实验全为单臂(UR5/Franka), "双臂/bimanual"任务仅见其 IJRR 2024 扩展版; 此处"双臂动作"归属该 RSS 2023 文有误, 键/会议无误} \\
leader--follower 动作分解 & 2024-25 & 把双臂动作预测分解为“leader 臂预测 $+$ follower 条件于 leader”两步——\emph{和 \cref{sec:ma-arch} 的 leader--follower 架构同构}，只是从控制律变成了网络结构 \\
$\pi_0$ 等 VLA & 2024-25 & 视觉-语言-动作大模型直出双臂 $+$ 移动控制 \\
\bottomrule
\end{tabular}
\end{table}

\begin{insight}[学习范式把经典概念从“控制律”变成“归纳偏置/隐式表征”]
学习范式没有抛弃本章的概念，而是把它们\emph{从“人工设计的控制律”变成了“网络的归纳偏置或隐式表征”}。最典型的是 leader--follower 动作分解：经典控制里它是 \cref{sec:ma-arch} 的架构选择（leader 持意图、follower 一致性跟随），在学习里它变成网络结构（follower 分支条件于 leader 分支的预测）——同一个“主从/leader--follower”思想，换了实现载体。理解经典的对称/相对/内力概念，反而能帮你看懂学习方法在隐式地处理什么、它们的归纳偏置合不合理。
\end{insight}

\begin{table}[H]\centering\small
\caption{经典控制（本章）vs 学习范式：互补而非替代}
\label{tab:ma-classic-vs-learning}
\begin{tabular}{lll}
\toprule
维度 & 经典控制（本章） & 学习范式 \\
\midrule
内力 & 显式建模、可精确调控 & 隐式（位置控制为主，力 demo 难采）——\emph{当前短板} \\
可分析 & 强（稳定性、对偶可证） & 弱（黑盒） \\
泛化到新任务 & 需重新设计 & 数据内泛化好、外推弱 \\
推广到多机/$N$ 臂 & 自然（本章主线） & 困难（多 agent 数据爆炸） \\
精细/接触任务 & 强（内力控制） & 精度高但力控弱 \\
\bottomrule
\end{tabular}
\end{table}

\begin{insight}[显式内力控制不会被学习取代]
ALOHA 体系主要做\emph{位置控制}，力控 demo 数据难采集——这正是 \cref{sec:ma-internalforce} 的显式内力控制不会被学习取代的原因：对内力敏感的精细装配、易碎物操作，仍需要本章的内力建模。一个活跃方向是把经典内力控制作为“安全外壳”包裹学习策略——学习策略给运动意图，经典内力环守护内力安全。这是经典控制与学习范式融合的典型形态，也呼应“MARL $+$ 传统规控混合”的主题。
\end{insight}

\begin{pitfall}[概念误区：认为“有了 ALOHA 这类学习方法，经典双臂控制就过时了”]
新手想法：“直接学就行，何必推导相对雅可比、内力分解这些复杂数学？”实际上：(1) 学习范式当前主要做位置控制，对内力敏感的任务（精密装配、易碎物）力控仍是短板，需经典内力控制兜底；(2) 学习方法难以推广到多臂/$N$ 机（数据爆炸），而本章的闭链数学天然推广；(3) 学习方法是黑盒，安全攸关场景需要经典控制的可分析性。正确理解：经典控制与学习范式互补——学运动意图、经典守内力与安全；理解经典概念还能帮你看懂学习方法在隐式处理什么。
\end{pitfall}

\begin{exercise}
\begin{enumerate}
  \item \textbf{[定位]} 给出 $5$ 个具体工程任务（如“双臂装配齿轮箱”“两台移动机械臂抬门”“学习双臂叠衣服”），对每个判断应主看本章还是单系统双臂内容，说明判据。
  \item \textbf{[分析]} ALOHA 的遥操作采数据是主从结构（人控 leader 臂）。论证：为什么遥操作主从在采数据时可行（人脑当内力环），但训练出的策略部署时若仍是纯位置控制，在内力敏感任务上会失败？如何用 \cref{sec:ma-internalforce} 的内力环改进？
  \item \textbf{[思考·研究级]} 设计“把经典内力控制作为学习策略的安全外壳”的具体方案：学习策略输出物体运动意图（绝对通道），经典内力环（\cref{sec:ma-internalforce}）守护内力安全（相对通道）。说明两者如何在 \cref{sec:ma-coopkin} 的绝对/相对分解下自然结合。
\end{enumerate}
\end{exercise}

% ════════════════════════════════════════════════════════════════
\section*{本章常见误解汇总}
\addcontentsline{toc}{section}{本章常见误解汇总}

把全章散落的“概念误区/思维陷阱”收口成一张表，便于回查（\cref{tab:ma-misconceptions}）。

\begin{table}[H]\centering\small
\caption{双臂/多臂协同操作常见误解辨正}
\label{tab:ma-misconceptions}
\begin{tabular}{p{5.6cm}p{7.0cm}p{1.4cm}}
\toprule
误解 & 正确理解 & 出处 \\
\midrule
多机协作视角只是换说法，对实现无影响 & 视角决定模块边界 $\to$ 代码结构、容错、可扩展性；视角 (A) 加臂要重写，视角 (B) 增量扩展 & \cref{sec:ma-perspective} \\
协同操作 $=$ 协同搬运（已讲，本章重复） & 搬运关注力分配（常点接触抽象）；操作多一层运动学协调（相对雅可比、在手重定向、内力精调） & \cref{sec:ma-perspective} \\
独立控制两只手，只要目标相容就行 & 两个独立控制器的跟踪误差在相对位姿上累积 $=$ 内力爆炸；必须用绝对/相对解耦后分别做运动/力控 & \cref{sec:ma-coopkin} \\
相对雅可比就是 $\mathbf{J}_R-\mathbf{J}_L$ & 必须先用平移算子统一参考点再做差；本质是 $\SE(3)$ 相对位姿的微分，不是矩阵减法 & \cref{sec:ma-coopkin} \\
有了 $\mathbf{J}_{\text{rel}}$ 投影就够，$\mathbf{J}_{\text{abs}}$ 没用 & $\mathbf{J}_{\text{abs}}$ 管“物体往哪走”，缺了它能不挤压物体却无法精确控制物体位置；内力解耦也需要它 & \cref{sec:ma-coopkin} \\
内力是浪费，应调成零 & 内力是抓取稳定性来源（防滑），目标是“防滑下界 $\leftrightarrow$ 强度上界”的区间，不是零 & \cref{sec:ma-duality,sec:ma-internalforce} \\
$\ker(\mathbf{G})$ 是固定子空间 & $\mathbf{G}$ 依赖 $\mathbf{r}_i$，物体一动 $\ker(\mathbf{G})$ 就旋转；开环内力会泄漏到运动通道 & \cref{sec:ma-duality} \\
内力可以开环施加 & 内力子空间随构型旋转，开环力会“漏”到运动通道；必须用相对位姿偏差闭环 & \cref{sec:ma-duality,sec:ma-internalforce} \\
内力阻抗刚度越大越稳 & $\mathbf{K}_i$ 过大 $\to$ 微小误差产生巨大内力 $\to$ 压坏物体/饱和，且与外环柔顺矛盾 & \cref{sec:ma-internalforce} \\
$N$ 臂内力 $=N-1$ 组双臂内力拼接 & $N$ 臂内力是整体 $6N-6$ 维子空间，配对只是一组（非唯一）基，且依赖末端运动学 & \cref{sec:ma-internalforce} \\
主从更简单所以更可靠 & 主从把内力管理“扫到地毯下”（从臂跟踪误差副产品），内力敏感任务反而最不可靠 & \cref{sec:ma-arch} \\
leader--follower 一致性 $=$ 经典主从 & 经典主从硬复现 leader 位姿（单向、放大抖动）；一致性是协商收敛（多跳、滤波、去中心化） & \cref{sec:ma-arch} \\
非对称任务 $=$ 主从控制 & 非对称是任务属性（$\alpha$），主从是架构属性；可用对称架构执行非对称任务且内力显式 & \cref{sec:ma-symmetry} \\
对称表述（绝对 $=$ 中点）适用所有任务 & 对称表述对对称任务最优；非对称任务“中点”无清晰物理意义，应用 $\alpha$ 以参考手为绝对变量 & \cref{sec:ma-symmetry} \\
$N$ 臂力分配照搬双臂“平分” & $N\ge3$ 有 $6N-6$ 维内力自由度待优化；异构/偏置下平分让弱臂先饱和，须用 QP & \cref{sec:ma-allocation} \\
分布式 ADMM 一定比集中 QP 慢 & 取决于 $N$、是否分布式硬件、通信代价；$N$ 大分布式部署时 ADMM 可能更快且无单点 & \cref{sec:ma-allocation} \\
有了 ALOHA 学习方法，经典控制过时 & 学习当前主要做位置控制、力控是短板、难推广多机、黑盒；与经典互补 & \cref{sec:ma-boundary} \\
\bottomrule
\end{tabular}
\end{table}

% ════════════════════════════════════════════════════════════════
\section*{本章小结}
\addcontentsline{toc}{section}{本章小结}

本章站在\emph{多机协作}的立场重讲了双臂与多臂协同操作。一句话总括认知主线：双臂是“$N=2$ 的多机协同操作”——它在最少自由度上同时具备多机系统的全部核心矛盾（共享目标一致、资源/力分配、异构容错），而把“臂”换成“腿”或加上浮动基座、把 $N$ 从 $2$ 加到任意，同一套闭链数学（相对雅可比、grasp matrix、内力分配、绝对/相对解耦）原样复用；协同操作的“集中 vs 分布”权衡（主从 vs 协调、集中 QP vs ADMM）与多机主题完全同构。

\textbf{视角重置（\cref{sec:ma-perspective}）}：把双臂当 $N=2$ 多机协同而非一台 $14$ 自由度机器——视角决定模块边界、容错与可扩展性，并能复用共识、任务分配、分布式优化、grasp matrix 全部工具。

\textbf{协同运动学（\cref{sec:ma-coopkin}）}：Uchiyama--Dauchez 对称表述把双臂运动分解为绝对（$\mathbf{J}_{\text{abs}}=\tfrac12(\bar{\mathbf{J}}_L+\bar{\mathbf{J}}_R)$，物体运动）$+$ 相对（$\mathbf{J}_{\text{rel}}=\bar{\mathbf{J}}_R-\bar{\mathbf{J}}_L$，内力通道）两个解耦通道（\cref{eq:ma-Jabs,eq:ma-Jrel}），$\alpha$ 把对称/非对称统一进一个连续参数。

\textbf{约束的两个面孔（\cref{sec:ma-duality}）}：虚功原理给出 $\delta\bar{\boldsymbol{\mathcal{V}}}_{\text{c}}=\mathbf{G}^\top\delta\boldsymbol{\mathcal{V}}_{\text{obj}}$（\cref{eq:ma-Gtransp}），证明速度约束 $\ker(\mathbf{J}_{\text{rel}})$ 与力约束 $\ker(\mathbf{G})$ 是同一闭链的对偶面孔——内力（$\ker(\mathbf{G})$）是相对运动（$\ker(\mathbf{G}^\top)$）的力共轭、必须闭环。

\textbf{内力控制（\cref{sec:ma-internalforce}）}：$\mathbf{f}_{\text{int}}=\tfrac12(\mathbf{h}_R-\mathbf{h}_L)$（\cref{eq:ma-fint}）落在 $\ker(\mathbf{G})$；Caccavale 内/外阻抗（外环管绝对、内环管内力）是对偶理论的控制器化身（\cref{eq:ma-tau-decomp}）；$N$ 臂推广为 $\mathbf{h}=\mathbf{G}^{+}\mathbf{F}+\mathbf{W}\boldsymbol{\lambda}$（\cref{eq:ma-narm-param}），多臂内力依赖末端运动学（Erhart--Hirche）。

\textbf{架构辩论（\cref{sec:ma-arch}）}：主从（集中式，内力失控）vs 对称协调（分布式，内力显式）vs leader--follower 一致性（去中心化，主从的软化，\cref{eq:ma-leaderfollow}），对应集中/分布/去中心化三分。

\textbf{对称/非对称（\cref{sec:ma-symmetry}）}：任务维度（$\alpha$）$\perp$ 架构维度；扩展相对雅可比处理非对称任务，非对称内力常是单向接触力；角色分配是任务分配的退化形式。

\textbf{多臂分配（\cref{sec:ma-allocation}）}：$N\ge3$ 力分配是带摩擦锥/能力约束的 QP（\cref{eq:ma-qp}），应最小化“最大归一化负载”；分布式版用 ADMM 把 QP 拆给各 agent，$\boldsymbol{\lambda}$ 是“物体 wrench 的协调价格”。

\textbf{边界与学习（\cref{sec:ma-boundary}）}：与单系统双臂内容互补；经典控制（学意图 $+$ 守内力）与学习范式互补，显式内力控制不会被学习取代。

\subsection*{符号表}
\begin{table}[H]\centering\small
\caption{本章主要符号}
\label{tab:ma-symbols}
\begin{tabular}{ll}
\toprule
符号 & 含义 \\
\midrule
$N$ & 协同操作的臂/接触点数目（双臂 $N=2$） \\
$\mathbf{q},\dot{\mathbf{q}}$ & 联合关节向量 / 关节速度（双 $7$ 自由度时 $\in\mathbb{R}^{14}$） \\
$\boldsymbol{\mathcal{V}}_i=[\mathbf{v}_i;\boldsymbol{\omega}_i]$ & 第 $i$ 臂末端 twist（本章线在前、角在后） \\
$\mathbf{X}(\mathbf{d})$ & twist 平移算子 $\big[\begin{smallmatrix}\mathbf{I}&-[\mathbf{d}]_\times\\\mathbf{0}&\mathbf{I}\end{smallmatrix}\big]$（统一参考点） \\
$\bar{\boldsymbol{\mathcal{V}}}_i,\bar{\mathbf{J}}_i$ & 统一到公共参考点 $\mathbf{p}_c$ 后的 twist / 雅可比 \\
$\mathbf{J}_{\text{abs}},\mathbf{J}_{\text{rel}}$ & 绝对雅可比 $\tfrac12(\bar{\mathbf{J}}_L+\bar{\mathbf{J}}_R)$ / 相对雅可比 $\bar{\mathbf{J}}_R-\bar{\mathbf{J}}_L$ \\
$\mathbf{J}_{\text{coop}},\mathbf{J}_{\text{aug}},\mathbf{J}_{\text{ext}}$ & 协同 / 增广 / 扩展（非对称）雅可比 \\
$\mathbf{T}_{\text{abs}},\mathbf{T}_{\text{rel}}$ & 绝对（物体）位姿 / 相对位姿 $\in\SE(3)$ \\
$\alpha$ & 绝对变量加权参数（$0.5$ 对称、$1/0$ 非对称参考手） \\
$\mathbf{G}$ & grasp matrix，$\mathbf{F}_{\text{load}}=\mathbf{G}\mathbf{h}$，$\in\mathbb{R}^{6\times6N}$（含力矩） \\
$\mathbf{h}=[\mathbf{h}_1;\dots;\mathbf{h}_N]$ & 各末端施加到物体的 wrench 堆叠 \\
$\mathbf{F}_{\text{load}}$ & 物体净 wrench（运动力） \\
$\mathbf{f}_{\text{int}}$ & 内力（wrench），落在 $\ker(\mathbf{G})$ \\
$\mathbf{W}$ & 内力子空间基矩阵（$\operatorname{Im}(\mathbf{W})=\ker(\mathbf{G})$） \\
$\boldsymbol{\lambda}$ & 内力参数 / ADMM 对偶变量（物体 wrench 协调价格） \\
$\mathbf{M}_o,\mathbf{D}_o,\mathbf{K}_o$ & 外环 object impedance 惯量/阻尼/刚度 \\
$\mathbf{M}_i,\mathbf{D}_i,\mathbf{K}_i$ & 内环 internal impedance 惯量/阻尼/刚度 \\
$\mu,f_{\min},f_{\text{crush}}$ & 摩擦系数 / 最小法向力 / 物体压坏阈值 \\
$\mathbf{L},\lambda_2$ & 通信图 Laplacian / 代数连通性（\cref{ch:mr-consensus}） \\
$\rho$ & ADMM 罚参数 \\
\bottomrule
\end{tabular}
\end{table}

\subsection*{定理/结论速查表}
\begin{table}[H]\centering\small
\caption{本章核心结论速查}
\label{tab:ma-keyresults}
\begin{tabular}{p{3.4cm}p{7.2cm}p{2.6cm}}
\toprule
结论 & 内容 & 出处 \\
\midrule
相对/绝对雅可比 & 统一参考点后 $\mathbf{J}_{\text{rel}}=\bar{\mathbf{J}}_R-\bar{\mathbf{J}}_L$、$\mathbf{J}_{\text{abs}}=\tfrac12(\bar{\mathbf{J}}_L+\bar{\mathbf{J}}_R)$ & \cref{eq:ma-Jabs,eq:ma-Jrel} \\
运动学-静力学对偶 & 虚功 $\Rightarrow\delta\bar{\boldsymbol{\mathcal{V}}}_{\text{c}}=\mathbf{G}^\top\delta\boldsymbol{\mathcal{V}}_{\text{obj}}$；$\mathbf{G}$ 力、$\mathbf{G}^\top$ 速度 & \cref{eq:ma-Gtransp} \\
两零空间对偶\rebuilt{待修: 原表述"内力 $\ker(\mathbf{G})\leftrightarrow$ 相对运动 $\ker(\mathbf{G}^\top)$ 同维同构"有误; 正确: 相对运动空间为 $\operatorname{Im}(\mathbf{G}^\top)$ 的正交补(非 $\ker(\mathbf{G}^\top)$), 见 \cref{sec:ma-nullspaces}} & 内力 $\ker(\mathbf{G})$（维 $6N-6$）$\leftrightarrow$ 破坏约束的相对运动（$\operatorname{Im}(\mathbf{G}^\top)$ 正交补，同维） & \cref{eq:ma-final-equiv} \\
双臂内力分解 & $\mathbf{f}_{\text{int}}=\tfrac12(\mathbf{h}_R-\mathbf{h}_L)$，$\mathbf{h}_{L/R}=\tfrac12\mathbf{F}\mp\mathbf{f}_{\text{int}}$ & \cref{eq:ma-fint,eq:ma-hLR} \\
Caccavale 双层 & 外环 object（绝对）$+$ 内环 internal（相对），正交解耦 & \cref{eq:ma-objimp,eq:ma-intimp,eq:ma-tau-decomp} \\
$N$ 臂内力参数化 & $\mathbf{h}=\mathbf{G}^{+}\mathbf{F}+\mathbf{W}\boldsymbol{\lambda}$，$\operatorname{Im}(\mathbf{W})=\ker(\mathbf{G})$ & \cref{eq:ma-narm-param} \\
leader--follower 一致性 & 带 leader 牵引项的共识，去中心化软化主从 & \cref{eq:ma-leaderfollow} \\
多臂 QP 分配 & $\min\tfrac12\mathbf{h}^\top\mathbf{Q}\mathbf{h}$ s.t. $\mathbf{G}\mathbf{h}=\mathbf{F}$、摩擦锥、力矩上限 & \cref{eq:ma-qp} \\
ADMM 分布式分配 & 各 agent 小 QP $+$ $\boldsymbol{\lambda}$（物体 wrench 价格）协调 & \cref{eq:ma-admm-local} \\
\bottomrule
\end{tabular}
\end{table}

\subsection*{知识点总表}
\begin{table}[H]\centering\small
\caption{本章知识点与难度}
\label{tab:ma-knowledge}
\begin{tabular}{p{3.6cm}p{7.6cm}c}
\toprule
知识点 & 核心要点 & 难度 \\
\midrule
多机视角看双臂 & $N=2$ 特例；翻译表；耦合强度谱（通信$\to$碰撞$\to$力$\to$力$+$运动） & $\bigstar\bigstar$ \\
对称表述 & 绝对 $=$ 中点、相对 $=$ 右减左；二体问题质心-相对坐标类比 & $\bigstar\bigstar\bigstar$ \\
相对/绝对雅可比 & 统一参考点 $\mathbf{X}(\mathbf{p}_c-\mathbf{p}_i)$ 后做和/差 & $\bigstar\bigstar\bigstar$ \\
协调逆运动学 & 绝对通道跟踪物体 $+$ 相对通道锁零；分层式优先相对约束 & $\bigstar\bigstar\bigstar$ \\
$\alpha$ 统一对称/非对称 & 加权绝对变量；任务对称性的连续参数 & $\bigstar\bigstar\bigstar$ \\
运动学-静力学对偶 & 虚功原理 $\to\mathbf{G}^\top$ 分配速度；$\mathbf{G}$ 力、$\mathbf{G}^\top$ 速度 & $\bigstar\bigstar\bigstar$ \\
$\ker(\mathbf{G})/\ker(\mathbf{J}_{\text{rel}})$ 对偶 & 内力 $\leftrightarrow$ 相对运动（同维）\pz{待修: 原"$\ker(\mathbf{G}^\top)$/SVD 配对"叙述误, 见 \cref{sec:ma-nullspaces}} & $\bigstar\bigstar\bigstar$ \\
内力必须闭环 & $\ker(\mathbf{G})$ 随构型旋转，开环力泄漏到运动通道 & $\bigstar\bigstar\bigstar$ \\
双臂内力分解 & $\mathbf{f}_{\text{int}}=\tfrac12(\mathbf{h}_R-\mathbf{h}_L)$；防滑下界 $\leftrightarrow$ 强度上界区间 & $\bigstar\bigstar\bigstar$ \\
Caccavale 内/外阻抗 & 外环 object（绝对）$+$ 内环 internal（相对）；正交解耦 & $\bigstar\bigstar\bigstar$ \\
$N$ 臂内力推广 & $\mathbf{h}=\mathbf{G}^{+}\mathbf{F}+\mathbf{W}\boldsymbol{\lambda}$；多臂内力依赖运动学 & $\bigstar\bigstar\bigstar\bigstar$ \\
主从 vs 协调 & 对应集中/分布；主从内力失控；协调内力显式 & $\bigstar\bigstar\bigstar$ \\
leader--follower 一致性 & 共识 $+$ leader 牵引项；去中心化软化主从 & $\bigstar\bigstar\bigstar$ \\
对称/非对称协调 & 任务维度 $\perp$ 架构维度；扩展相对雅可比；非对称内力常单向 & $\bigstar\bigstar\bigstar$ \\
多臂 QP 力分配 & $\min\mathbf{h}^\top\mathbf{Q}\mathbf{h}$ s.t. $\mathbf{G}\mathbf{h}=\mathbf{F}$、摩擦锥、力矩上限 & $\bigstar\bigstar\bigstar\bigstar$ \\
ADMM 分布式分配 & 各 agent 小 QP $+$ $\boldsymbol{\lambda}$（物体 wrench 价格）协调 & $\bigstar\bigstar\bigstar\bigstar$ \\
两套双臂内容分工 & 单系统（工程深度）vs 多机（结构推广），互补 & $\bigstar\bigstar$ \\
经典 vs 学习范式 & 学意图 $+$ 经典守内力；力控是学习短板 & $\bigstar\bigstar$ \\
\bottomrule
\end{tabular}
\end{table}

% ════════════════════════════════════════════════════════════════
\section*{延伸阅读}
\addcontentsline{toc}{section}{延伸阅读}
按“先读哪个、为什么读”组织。
\begin{itemize}
  \item \textbf{经典控制线（本章主线源头）}：Uchiyama--Dauchez 对称运动学表述与非主从协调控制\cite{uchiyama1993symmetric}（\cref{sec:ma-coopkin,sec:ma-arch,sec:ma-symmetry} 的理论源头，配套早期 ICRA 1988 hybrid position/force 方案\cite{uchiyama1987symmetric}\pz{存疑: 年份 1988 非 1987, 见 \cref{tab:ma-history}}）；Chiacchio 等多臂系统任务空间动力学分析\cite{chiacchio1996task}\pz{存疑: 该键的论文/卷期与"绝对/相对变量"内容归属可能张冠李戴, 详见 \cref{tab:ma-history} 标注}（绝对/相对变量现代形式）；Caccavale 等双臂 $6$ 自由度阻抗\cite{caccavale2008sixdof}（\cref{sec:ma-internalforce} 主线，内/外阻抗严格 $6$-DOF 刚度）；Jamisola--Roberts 紧凑相对雅可比\cite{jamisola2015compact}（\cref{sec:ma-compact}）；Erhart--Hirche 内力分析与负载分配\cite{erhart2017internal}\pz{存疑: 年份应为 2015 非 2017, 见 \cref{tab:ma-history}}（\cref{sec:ma-internalforce,sec:ma-allocation} 多臂内力严格基础，$N\ge3$ 必读）。
  \item \textbf{多机/分布式线}：Khatib augmented object（object-level 动力学）\cite{khatib1988object}\rebuilt{原误用键 khatib1987operational(1987 操作空间原文, 不含 augmented object), 已改 khatib1988object(ISRR 1988), 见 \cref{tab:ma-history}}（运动/力解耦的动力学基础）；Verginis--Dimarogonas 内力调节的刚性理论视角\cite{verginis2019cooperative}\pz{存疑: arXiv 应为 1911.01297、应含第三作者 Zelazo、期刊 T-CNS, 见 \cref{tab:ma-history}}（\cref{sec:ma-internalforce,sec:ma-allocation} 分布式内力调节，无需负载惯性）；Almeida--Karayiannidis 等扩展相对雅可比\cite{almeida2019asymmetric}\pz{存疑: 第二作者 Karayiannidis 非"Lima"、会议 ISRR 2019 非 IROS, 见 \cref{tab:ma-history}}（\cref{sec:ma-symmetry} 非对称任务）。
  \item \textbf{学习范式线}：Zhao--Finn 等 ALOHA\cite{zhao2023aloha}（双臂遥操作 $+$ ACT 奠基）；Chi 等 Diffusion Policy\cite{chi2023diffusion}（DDPM 生成动作 chunk）\rebuilt{待修: 原文 RSS 2023 实验为单臂, "双臂"任务仅见其 IJRR 2024 扩展版, 见 \cref{tab:ma-learning} 标注}。
  \item \textbf{教材/综述}：Springer Handbook of Robotics 的 “Cooperative Manipulation” 章（Caccavale \& Uchiyama 执笔）是本章经典控制内容的权威综述与符号标准。\pz{该手册章节、Siciliano 等《Robotics: Modelling, Planning and Control》（雅可比/伪逆/阻抗前置）尚未加入 \texttt{refs.bib}，待补键。}
\end{itemize}

\section*{本章与后续章节的关系}
\addcontentsline{toc}{section}{本章与后续章节的关系}
本章是协同操作的“闭链母板”，为后续提供数学基础并向多智能体学习输送概念（\cref{tab:ma-relations}）。\pz{下表“后续章节”指代多足 loco-manipulation、人形全身协同、MARL 等章；这些章尚未写就，编号留待确定后改 \texttt{\textbackslash cref}。}

\begin{table}[H]\centering\small
\caption{本章与后续章节的关系}
\label{tab:ma-relations}
\begin{tabular}{p{3.4cm}p{5.6cm}p{4.4cm}}
\toprule
后续章节 & 关系 & 本章铺垫的知识点 \\
\midrule
多足协同 loco-manipulation & \emph{直接复用}——“臂”换“腿”，每条腿是“接触点 $+$ 运动学”，闭链约束、grasp matrix、内力分配原样搬过去，只是雅可比多浮动基座列、约束多接触维持 & 相对/绝对雅可比、对偶与内力闭环、力分配 QP/ADMM \\
人形/四足 $+$ 臂全身协同 & \emph{扩展}——双臂装到浮动基座，绝对/相对分解 $+$ WBC 分层；多肢体协同是 $N$ 接触点的统一闭链 & 对称表述、内/外阻抗（升级为全身）、$N$ 接触点内力 \\
多智能体强化学习 & \emph{对照与融合}——经典控制作 MARL 安全外壳（学意图 $+$ 经典守内力）；leader--follower 在 MARL 里变 CTDE 角色结构 & 主从/一致性架构、经典 vs 学习互补 \\
MARL $+$ 传统规控混合 & \emph{核心呼应}——“经典内力环 $+$ 学习运动策略”正是本章对该主题的贡献 & 内力闭环作安全约束、混合范式 \\
\bottomrule
\end{tabular}
\end{table}

\noindent\textbf{前向预告}：多足 loco-manipulation 的“足端力分配”和本章 \cref{sec:ma-allocation} 的“多臂负载分配”是同一个 QP——只是接触点是脚而非手，且多了“摆动腿不出力”的接触序列约束。本章的闭链数学不是双臂专属，而是“任意一组末端共持负载”的通用语言——这正是把它放在“多机协作”而非“机械臂”方向的根本理由。

\section*{故障排查手册}
\addcontentsline{toc}{section}{故障排查手册}
协同操作调试时最常见的 $5$ 类故障（\cref{tab:ma-troubleshoot}）。

\begin{table}[H]\centering\small
\caption{协同操作故障排查（症状 $\to$ 原因 $\to$ 排查 $\to$ 相关节）}
\label{tab:ma-troubleshoot}
\begin{tabular}{p{3.0cm}p{4.6cm}p{5.2cm}}
\toprule
症状 & 可能原因 & 排查步骤（含相关节） \\
\midrule
物体被无故挤压 / 两臂“对抗”（内力失控） & (a) $\mathbf{J}_{\text{rel}}$ 用 \texttt{LOCAL\_WORLD\_ALIGNED} 直接 $\mathbf{J}_R-\mathbf{J}_L$ 没统一参考点；(b) 内力开环；(c) 两臂时钟不同步 & 构造物体纯转动 $\dot{\mathbf{q}}$ 测 $\|\mathbf{J}_{\text{rel}}\dot{\mathbf{q}}\|$（\cref{sec:ma-coopkin}）；查内力是否 internal impedance 闭环（\cref{sec:ma-internalforce}）；查 $\mathbf{x}_{\text{abs}}^{\text{des}}(t)$ 是否同一时钟（\cref{sec:ma-arch}） \\
物体能搬动但轨迹漂移 / 内力随姿态缓慢增长 & (a) 内力开环 $\to\ker(\mathbf{G})$ 旋转、力泄漏到运动通道；(b) 内外环 wrench 在不同参考点叠加；(c) 缺重力/惯性补偿 & 把内力改成相对位姿偏差闭环看漂移是否消失；确认 $\mathbf{J}_{\text{abs}},\mathbf{J}_{\text{rel}}$ 与内外环 wrench 同一参考点（\cref{sec:ma-duality,sec:ma-internalforce}）；加重力补偿前馈 \\
主臂经过某构型时内力剧烈震荡（主从架构） & 主臂接近奇异 $\to$ 实际运动滞后/抖动 $\to$ 从臂复现抖动 $\to$ 相对位姿剧烈波动 & 监测主臂操作度 $w=\sqrt{\det(\mathbf{J}_L\mathbf{J}_L^\top)}$ 确认震荡是否在 $w$ 低谷；换对称协调或 leader--follower 一致性（\cref{sec:ma-arch}） \\
$N\ge3$ 臂搬运时某臂率先力矩饱和、连锁过载 & (a) 朴素平分而非优化分配；(b) QP 代价没编码能力差异；(c) 没设最小法向力 & 改用 QP 分配（加权或 min--max 归一化力）；按 $1/\tau_{\max,i}$ 加权；加 $f_{\min}$ 约束防滑（\cref{sec:ma-allocation}） \\
摩擦锥约束设了但接触仍打滑 & (a) 二阶锥直接喂给只支持线性约束的 QP 求解器；(b) $\mu$ 估计过乐观；(c) 接触法向建模错误 & 摩擦锥用金字塔线性化后喂 QP 或换 SOCP 求解器；保守取 $\mu$；核对每个接触法向方向（\cref{sec:ma-allocation}） \\
\bottomrule
\end{tabular}
\end{table}
